\textbf{Examples}\par

\begin{enumerate}
\def\labelenumi{(\arabic{enumi})}
\item
  For every \(a \in \mathbf{A}\) , the set of \(a^n\) , where \(n \in \mathbf{N}\) , is a multiplicative subset of \(\mathbf{A}\) .
\item
  Let \(\mathfrak{p}\) be an ideal of \(\mathbf{A}\) . For \(\mathbf{A} - \mathfrak{p}\) to be a multiplicative subset of \(\mathbf{A}\) , it is necessary and sufficient that \(\mathfrak{p}\) be prime.
\item
  The set of elements of A which are not divisors of zero is a multiplicative subset of A.
\item
  If S and T are multiplicative subsets of A, the set ST of products st, where \(s \in S\) and \(t \in T\) , is a multiplicative subset.
\item
  Let \(\mathfrak{S}\) be a directed set (with respect to the relation \(\subset\) ) of multiplicative subsets of A. Then \(T = \bigcup_{S \in \mathfrak{S}} S\) is a multiplicative subset of A, as any two elements of T belong to some subset \(S \in \mathfrak{S}\) , hence their product belongs to T.
\item
  Every intersection of multiplicative subsets of A is a multiplicative subset.
\end{enumerate}

For every subset S of a ring A, there exist multiplicative subsets of A containing S, for example A itself. The intersection of all these subsets is the smallest multiplicative subset of A containing S; it is said to be generated by S. It follows immediately that it is the set consisting of all the finite products of elements of S.

PROPOSITION 1. Let A be a ring and S a subset of A. There exists a ring A' and a homomorphism h of A to A' with the following properties:

\begin{enumerate}
\def\labelenumi{(\arabic{enumi})}
\item
  the elements of \(h(\mathbf{S})\) are invertible in \(\mathbf{A}'\) ;
\item
  for every homomorphism \(u\) of \(\mathbf{A}\) to a ring \(\mathbf{B}\) such that the elements of \(u(\mathbf{S})\) are invertible in \(\mathbf{B}\) , there exists a unique homomorphism \(u'\) of \(\mathbf{A}'\) to \(\mathbf{B}\) such that \(u = u' \circ h\) .
\end{enumerate}

In other words, \((A', h)\) is a solution of the universal mapping problem (Set Theory, Chapter IV, §3, no. 1) with the following conditions: the species of structure \(\Sigma\) considered is that of a ring, the morphisms are ring homomorphisms and the \(\alpha\) -mappings are homomorphisms of A to a ring such that the image of S under such a homomorphism consists of invertible elements. Recall (loc. cit.) that, if \((A', h)\) and \((A_{1}', h_{1})\) are both solutions of this problem, there exists a unique isomorphism \(j: A' \to A_{1}'\) such that \(h_{1} = j \circ h\) .

Let \(\bar{S}\) be the multiplicative subset of A generated by S. Clearly every solution of the above universal mapping problem is also a solution of the universal mapping problem obtained by replacing S by \(\bar{S}\) and conversely.

Consider, in the set \(\mathbf{A} \times \overline{\mathbf{S}}\) , the following relation between elements \((a, s)\) , \((a', s')\) :

\begin{enumerate}
\def\labelenumi{(\arabic{enumi})}
\tightlist
\item
 ``There exists \(t \in \overline{\mathbf{S}}\) such that \(t(sa' - s'a) = 0\)''.
\end{enumerate}

This relation is an equivalence relation: it is clearly reflexive and symmetric; it is transitive, for if \(t(sa' - s'a) = 0\) and \(t'(s'a'' - s''a') = 0\) , then \(tt's'(sa' - s'a) = 0\) and \(tt's' \in \overline{\mathbb{S}}\) . Let \(A'\) be the quotient set of \(A \times \overline{S}\) under this equivalence relation; for every ordered pair \((a, s) \in A \times \overline{S}\) we denote by \(a/s\) the canonical image of \((a, s)\) in \(A'\) and we set \(h(a) = a/1\) for all \(a \in A\) . We shall see that \(A'\) can be given a ring structure such that the ordered pair \((A', h)\) solves the problem.

Let \(x = a / s\) and \(y = b / t\) be two elements of \(A'\) . The elements \((ta + sb) / st\) and \(ab / st\) depend only on \(x\) and \(y\) ; for if \(x = a' / s'\) , there exists by hypothesis \(r \in \overline{\mathbb{S}}\) such that \(r(s'a - sa') = 0\) , whence

\[
r \left(s ^ {\prime} t (t a + s b) - s t \left(t a ^ {\prime} + s ^ {\prime} b\right)\right) = 0
\]

and \(r(s'tab - sta'b) = 0\) . It is easily verified that the laws of composition \((x,y) \mapsto x + y = (ta + sb)/st\) and \((x,y) \mapsto xy = ab/st\) define a commutative ring structure on \(A'\) , under which 0/1 is the identity element under addition and 1/1 is the unit element. Moreover, it is immediate that \(h\) is a ring homomorphism and that, for all \(s \in S\) , \(s/1\) is invertible in \(A'\) , its inverse being 1/s. Finally let B be a ring and \(u: A \to B\) a homomorphism such that the elements \(u(S)\) are invertible in B; there exists a unique mapping \(u': A' \to B\) such that

\[
u ^ {\prime} (a / s) = u (a) (u (s)) ^ {- 1} \quad (a \in A, s \in \bar {S}).\tag{2}
\]

If \(a / s = a' / s'\) , there exists \(t \in \overline{S}\) such that \(t(sa' - s'a) = 0\) , whence \(u(t)(u(s)u(a') - u(s'))u(a)) = 0\) and, as \(u(t), u(s)\) and \(u(s')\) are invertible, \(u(a)(u(s))^{-1} = u(a')(u(s'))^{-1}\) . It is easily verified that \(u'\) is a homomorphism with respect to addition and multiplication; finally, clearly \(u' \circ h = u\) and \(u'\) is the only homomorphism satisfying this relation, as it implies

\[
u ^ {\prime} (a / s) = u ^ {\prime} ((a / 1) (1 / s)) = u ^ {\prime} (1 / s) u ^ {\prime} (a / 1) = u ^ {\prime} (1 / s) u (a)
\]

and \(1 = u'(1/1) = u'(s/1)u'(1/s) = u(s)u'(1/s)\) , whence formula (2).

DEFINITION 2. Let A be a ring, S a subset of A and \(\overline{\mathbf{S}}\) the multiplicative subset generated by S. The ring of fractions of A defined by S and denoted by \(\mathrm{A}[\mathrm{S}^{-1}]\) is the quotient set of \(\mathrm{A} \times \overline{\mathrm{S}}\) under the equivalence relation (1) with the ring structure defined by

\[
(a / s) + (b / t) = (t a + s b) / s t, \quad (a / s) (b / t) = (a b) / (s t)
\]

for \(a, b\) in A, \(s, t\) in \(\bar{S}\) . The canonical mapping of A to A{[}S \(^{-1}\) {]} is the homomorphism \(a \mapsto a/1\) , which makes A{[}S \(^{-1}\) {]} into an A-algebra.

In this chapter we usually denote this canonical mapping by \(i_{A}^{S}\) ; the proof of Proposition 1 shows that the ordered pair \((\mathrm{A}[\mathrm{S}^{-1}], i_{\mathrm{A}}^{\mathrm{S}})\) satisfies the conditions of the statement of this proposition.

Remarks

\begin{enumerate}
\def\labelenumi{(\arabic{enumi})}
\item
  Clearly \(\mathbf{A}[\overline{\mathbf{S}}^{-1}] = \mathbf{A}[\mathbf{S}^{-1}]\) .
\item
  Two elements of A{[}S \(^{-1}\) {]} can always be written in the form a/s and a'/s (a, a' in A, s ∈ S) with the same ``denominator'' s, for if b/t and b'/t' are two elements of A{[}S \(^{-1}\) {]}, then b/t = bt'/tt' and b'/t' = b't/tt'.
\item
  The kernel of \(i_{\mathbf{A}}^{\mathbf{s}}\) is the set of \(a \in \mathbf{A}\) such that there exists \(s \in \overline{\mathbf{S}}\) satisfying \(sa = 0\) ; for \(i_{\mathbf{A}}^{\mathbf{s}}\) to be injective, it is necessary and sufficient that \(\mathbf{S}\) contain no divisor of zero in \(\mathbf{A}\) .
\item
  If \(S\) contains a nilpotent element, then \(0 \in \overline{S}\) and the ring \(A[S^{-1}]\) is reduced to 0; this follows easily from Definition 2.
\item
  For \(i_{\mathbf{A}}^{\mathbf{S}}\) to be a bijection, it is necessary and sufficient that every element \(s \in \mathbf{S}\) be invertible in A: the condition is obviously necessary, since \(s/1\) is invertible in \(\mathrm{A}[\mathrm{S}^{-1}]\) ; it is sufficient, since for all \(t \in \overline{\mathrm{S}}\) , \(t\) is therefore invertible in A and \(a/t = at^{-1}/1\) in \(\mathrm{A}[\mathrm{S}^{-1}]\) ; hence \(i_{\mathbf{A}}^{\mathbf{S}}\) is surjective and it has already been seen in Remark 3 that it is injective. Then A and \(\mathrm{A}[\mathrm{S}^{-1}]\) are identified by means of \(i_{\mathbf{A}}^{\mathbf{S}}\) .
\end{enumerate}

Example (7). If R is the set of elements in A which are not divisors of 0, the ring \(\mathrm{A}[\mathrm{R}^{-1}]\) is precisely what we have called the ring of fractions of A (Algebra, Chapter I, §9, no. 4); to avoid any confusion we shall often call it the total ring of fractions of A. In particular, if A is an integral domain, \(\mathrm{A}[\mathrm{R}^{-1}]\) is the field of fractions of A (loc. cit.).

PROPOSITION 2. Let A, B be two rings, S a subset of A, T a subset of B and f a homomorphism from A to B such that \(f(S) \subset T\) . There exists a unique homomorphism \(f'\) from A{[}S \(^{-1}\) {]} to B{[}T \(^{-1}\) {]} such that \(f'(a/1) = f(a)/1\) for all \(a \in A\) .

Suppose further that \(\mathbf{T}\) is contained in the multiplicative subset of \(\mathbf{B}\) generated by \(f(\mathbf{S})\) . Then, if \(f\) is surjective (resp. injective) so is \(f'\) .

The first assertion amounts to saying that there exists a unique homomorphism \(f'\) : \(A[S^{-1}] \to B[T^{-1}]\) giving a commutative diagram:

\[
\begin{array}{c} \mathrm{A} \xrightarrow {f} \mathrm{B} \\ i _ {\mathbf {A}} ^ {\mathbf {S}} \Biggl \downarrow \qquad \qquad \qquad \qquad \Biggl \downarrow i _ {\mathbf {B}} ^ {\mathbf {T}} \\ \mathrm{A} [ \mathrm{S} ^ {- 1} ] \xrightarrow [ f ^ {\prime} ]{} \mathrm{B} [ \mathrm{T} ^ {- 1} ] \end{array}
\]

Now the relation \(f(\mathbf{S}) \subset \mathbf{T}\) implies that \(i_{\mathbf{B}}^{\mathbf{T}}(f(s))\) is invertible in \(\mathbf{B}[\mathbf{T}^{-1}]\) for all \(s \in \mathbf{S}\) and it is sufficient to apply Proposition 1 to \(i_{\mathbf{B}}^{\mathbf{T}} \circ f\) . It follows easily from (2) that, for all \(a \in \mathbf{A}\) and \(s \in \overline{\mathbf{S}}\) (multiplicative subset of \(\mathbf{A}\) generated by \(\mathbf{S}\) ),

\[
f ^ {\prime} (a / s) = f (a) / f (s).\tag{3}
\]

Suppose that \(\mathbf{T}\) is contained in the multiplicative subset generated by \(f(\mathbb{S})\) , which is precisely \(f(\overline{\mathbb{S}})\) . Then it follows from (3) that, if \(f\) is surjective, so is \(f'\) .

Suppose now that \(f\) is injective. Let \(a / s\) be an element of the kernel of \(f'\) . As the multiplicative subset generated by \(T\) is \(f(\overline{S})\) , there is an element \(s_1 \in \overline{S}\) such that \(f(s_1)f(a) = 0\) , whence \(f(s_1a) = 0\) and consequently \(s_1a = 0\) since \(f\) is injective; then \(a / s = 0\) , which proves that \(f'\) is injective.

Remark (6). If the elements of T are invertible in B, \(B[T^{-1}]\) is identified with B by means of the isomorphism \(i_{B}^{T}\) and \(f'\) then becomes identical with the unique homomorphism \(u'\) of \(A[S^{-1}]\) to B such that \(u' \circ i_{A}^{S} = f\) .

COROLLARY 1. Let A be a ring, S a subset of A and u an injective homomorphism of A to a ring B such that the elements of \(u(S)\) are invertible in B. The unique homomorphism \(u'\) of \(\mathbf{A}[\mathbf{S}^{-1}]\) to B such that \(u' \circ i_{\mathbf{A}}^{\mathbf{S}} = u\) is then injective.

This is an immediate consequence of Proposition 2 and Remark 6.

COROLLARY 2. Let A be a ring and S and T two subsets of A such that \(\mathbf{S} \subset \mathbf{T}\) . There exists a unique homomorphism \(i_{\mathbf{A}}^{\mathrm{T},\mathrm{S}}\) from \(\mathrm{A}[\mathrm{S}^{-1}]\) to \(\mathrm{A}[\mathrm{T}^{-1}]\) such that \(i_{\mathbf{A}}^{\mathrm{T}} = i_{\mathbf{A}}^{\mathrm{T},\mathrm{S}} \circ i_{\mathbf{A}}^{\mathrm{S}}\) .

For all \(a \in \mathbf{A}\) , \(i_{\mathbf{A}}^{\mathbf{T},\mathbf{S}}\) then maps the element \(a / s\) in \(\mathbf{A}[\mathbf{S}^{-1}]\) to the element \(a / s\) in \(\mathbf{A}[\mathbf{T}^{-1}]\) .

Remark (7). Note that, if \(i_{A}^{T}\) is injective, so is \(i_{A}^{T,S}\) (Corollary 1). This is what happens if T is the set R of elements of A which are not divisors of 0; then it is possible to identify \(A[S^{-1}]\) with the subring of the total ring of fractions \(A[R^{-1}]\) generated by A and the inverses in \(A[R^{-1}]\) of the elements of S.

COROLLARY 3. Let A, B, C be three rings, S (resp. T, U) a multiplicative subset of A (resp. B, C), \(f\colon A \to B\) , \(g\colon B \to C\) two homomorphisms and \(h\colon A \to C\) the composite homomorphism \(g \circ f\) ; suppose that \(f(S) \subset T\) , \(g(T) \subset U\) . Let \(f'\colon A[S^{-1}] \to B[T^{-1}]\) , \(g'\colon B[T^{-1}] \to C[U^{-1}]\) , \(h'\colon A[S^{-1}] \to C[U^{-1}]\) the homomorphisms corresponding to \(f\) , \(g\) , \(h\) ; then \(h' = g' \circ f'\) .

This follows easily from the definitions.

In particular, if S, T, U are three multiplicative subsets of A such that \(S \subset T \subset U\) , then \(i_{A}^{U,S} = i_{A}^{U,T} \circ i_{A}^{T,S}\) .

COROLLARY 4. Let S be a subset of a ring A, B a subring of A{[}S \(^{-1}\) {]} containing \(i_{\text{A}}^{\text{S}}(\text{A})\) and S' the set \(i_{\text{A}}^{\text{S}}(\text{A})\) . Let j be the canonical injection of B into A{[}S \(^{-1}\) {]}; the unique homomorphism g from B{[}S \(^{-1}\) {]} to A{[}S \(^{-1}\) {]} such that \(g \circ i_{\text{A}}^{\text{S}} = j\) is an isomorphism.

The mapping g is injective by Corollary 1; the ring \(g(\mathbf{B}[\mathbf{S}^{\prime-1}])\) contains \(i_{\mathbf{A}}^{\mathbf{s}}(\mathbf{A})\) and the inverse of the elements of \(S'\) ; hence it is equal to \(A[S^{-1}]\) .

If A is an integral domain and \(0 \notin S\) , the notation \(A[S^{-1}]\) agrees with that of Algebra, Chapter IV, § 2, no. 1; also, if S is multiplicative, \(A[S^{-1}]\) coincides in this case with the set denoted by \(S^{-1}A\) in Algebra, Chapter I, § 1, no. 1.

As an extension of notation, for any multiplicative subset S of a ring A, we henceforth denote by \(S^{-1}A\) the ring of fractions \(A[S^{-1}]\) . If S is the complement of a prime ideal p of A, we write \(A_{p}\) instead of \(S^{-1}A\) .

If A is an integral domain and \(0 \notin S\) , \(S^{-1}A\) is always identified with a subring of the field of fractions of A, containing A (Remark 7).

\subsection{Modules of fractions}

The canonical homomorphism \(i_{\mathbf{A}}^{\mathbf{S}} \colon \mathbf{A} \to \mathbf{A}[\mathbf{S}^{-1}]\) defined in no. 1 allows us to consider every \(\mathbf{A}[\mathbf{S}^{-1}]\) -module as an A-module.

PROPOSITION 3. Let A be a ring, S a subset of A, M an A-module, M' the A-module M ⊗\_A A{[}S\^{}\{-1\}{]} and f the canonical A-homomorphism x ↦ x ⊗ 1 of M to M'. Then:

\begin{enumerate}
\def\labelenumi{(\arabic{enumi})}
\item
  For all \(s \in \mathbf{S}\) , the homothety \(z \mapsto sz\) of \(\mathbf{M}'\) is bijective.
\item
  For every A-module N such that, for all \(s \in S\) , the homothety \(y \mapsto sy\) of N is bijective, and every homomorphism u of M to N, there exists a unique homomorphism \(u'\) of \(M'\) to N such that \(u = u' \circ f\) .
\end{enumerate}

In other words, \((M',f)\) is a solution of the universal mapping problem (Set Theory, Chapter IV, §3, no. 1) with the following conditions: the species of structure \(\Sigma\) is that of an A-module in which the homotheties induced by the elements of S are bijective, the morphisms are A-module homomorphisms and the \(\alpha\) -mappings are also A-module homomorphisms.

For every A-module N and all \(a \in \mathbf{A}\) , denote by \(h_a\) the homothety \(y \mapsto ay\) in \(\mathbf{N}\) ; \(a \mapsto h_a\) is then a ring homomorphism from A to \(\operatorname{End}_{\mathbf{A}}(\mathbf{N})\) . To say that \(h_a\) is bijective means that \(h_a\) is an invertible element of \(\operatorname{End}_{\mathbf{A}}(\mathbf{N})\) . Suppose that, for all \(s \in \mathbf{S}\) , \(h_s\) is invertible in \(\operatorname{End}_{\mathbf{A}}(\mathbf{N})\) ; the elements \(h_a\) , where \(a \in \mathbf{A}\) , and the inverses of the elements \(h_s\) , where \(s \in \mathbf{S}\) , then generate in \(\operatorname{End}_{\mathbf{A}}(\mathbf{N})\) a commutative subring B and the homomorphism \(a \mapsto h_a\) from A to B is such that the images of the elements of S are invertible. Then it follows (no. 1, Proposition 1) that there exists a unique homomorphism \(h'\) of \(\mathrm{A}[\mathrm{S}^{-1}]\) to B such that

\[
h ^ {\prime} (a / s) = h _ {a} \left(h _ {s}\right) ^ {- 1};
\]

we know (Algebra, Chapter II, § 1, no. 14) that such a homomorphism defines on N an A{[}S \(^{-1}\) {]}-module structure such that (a/s).y = h \(_{s}^{-1}\) (a.y); the A-module structure derived from this A{[}S \(^{-1}\) {]}-module structure by means of the homomorphism i \(_{A}^{S}\) is precisely the structure given initially.

Conversely, if N is an A{[}S \(^{-1}\) {]}-module and it is considered as an A-module by means of \(i_{\text{A}}^{\text{S}}\) , the homotheties \(y \mapsto sy\) , for \(s \in \text{S}\) , are bijective, for \(y \mapsto (1/s)y\) is the inverse mapping of \(y \mapsto sy\) ; and the A{[}S \(^{-1}\) {]}-module structure on N derived from its A-module structure by the process described above is the A{[}S \(^{-1}\) {]}-module structure given initially. Thus there is a canonical one-to-one correspondence between A{[}S \(^{-1}\) {]}-modules and A-modules in which the homotheties induced by the elements of S are bijective; moreover, if N, N' are two A-modules with this property every A-module homomorphism \(u: \text{N} \to \text{N}'\) is also a homomorphism of the A{[}S \(^{-1}\) {]}-module structures of N and N', as, for all \(y \in \text{N}\) and all \(s \in \text{S}\) , we may write \(u(y) = u(s. ((1/s)y)) = s.u((1/s)y)\) , whence \(u((1/s)y) = (1/s)u(y)\) ; the converse is obvious.

This being so, the statement of Proposition 3 is just the characterization of the module obtained from M by extending the scalars to A{[}S \(^{-1}\) {]}, taking account of the above interpretation (Algebra, Chapter II, § 5, no. 1, Remark 1).

DEFINITION 3. Let A be a ring, S a subset of A, \(\overline{\mathbf{S}}\) the multiplicative subset of A generated by S and M an A-module. Then the module of fractions of M defined by S and denoted by \(\mathbf{M}[\mathbf{S}^{-1}]\) or \(\overline{\mathbf{S}}^{-1}\mathbf{M}\) is the \(\mathbf{A}[\mathbf{S}^{-1}]\) -module \(\mathbf{M} \otimes_{\mathbf{A}} \mathbf{A}[\mathbf{S}^{-1}]\) .

In this chapter we shall usually denote by \(i_{\mathbf{M}}^{\mathbf{S}}\) the canonical mapping \(m \mapsto m \otimes 1\) of \(\mathbf{M}\) to \(\mathbf{M}[\mathbf{S}^{-1}]\) .

Remarks

\begin{enumerate}
\def\labelenumi{(\arabic{enumi})}
\item
  Clearly \(\mathbf{M}[\overline{\mathbf{S}}^{-1}] = \mathbf{M}[\mathbf{S}^{-1}]\) .
\item
  For \(m \in \mathbf{M}\) and \(s \in \overline{\mathbf{S}}\) we also write \(m / s\) for the element \(m \otimes (1 / s)\) of \(\mathbf{M}[\mathbf{S}^{-1}]\) . Every element of \(\mathbf{M}[\mathbf{S}^{-1}]\) is of the form, for such an element is of the form \(\sum_{i} m_i \otimes (a_i / s)\) , where \(m_i \in \mathbf{M}\) , \(a_i \in \mathbf{A}\) , \(s \in \mathbf{S}\) (no. 1, Remark 2), and
\[
m _ {i} \otimes (a _ {i} / s) = (a _ {i} m _ {i}) \otimes (1 / s),
\]

hence \(\sum_{i} m_{i} \otimes (a_{i}/s) = m \otimes (1/s)\) , where \(m = \sum_{i} a_{i} m_{i}\) . Then


\[
(m / s) + (m ^ {\prime} / s ^ {\prime}) = (s ^ {\prime} m + s m ^ {\prime}) / s s ^ {\prime}\tag{4}
\]

\[
(a / s) (m / s ^ {\prime}) = (a m) / (s s ^ {\prime})\tag{5}
\]

where \(m, m' \in \mathbf{M}, a \in \mathbf{A}\) and \(s, s' \in \mathbf{S}\) .

\item
  If S is the complement of a prime ideal p of A, we write \(M_{p}\) instead of \(S^{-1}M\) .
\item
  Let M be an A{[}S \(^{-1}\) {]}-module; if M is considered canonically as an A-module, \(i_{M}^{S}\) is a bijection, for the ordered pair consisting of M and the identity mapping \(l_{M}\) is also trivially a solution of the universal mapping problem solved by M{[}S \(^{-1}\) {]} and \(i_{M}^{S}\) . Then M is identified with M{[}S \(^{-1}\) {]}.
\end{enumerate}

PROPOSITION 4. Let S be a multiplicative subset of A and M an A-module. For m/s = 0 (m ∈ M, s ∈ S), it is necessary and sufficient that there exist s' ∈ S such that s'm = 0.

If \(s' \in S\) is such that \(s'm = 0\) , clearly \(m/s = (s'm)/(ss') = 0\) . Conversely, suppose that \(m/s = 0\) . As \(1/s\) is invertible in \(S^{-1}A\) , \(m/1 = 0\) . For every sub-A-module \(P\) of \(S^{-1}A\) containing 1, we denote by \(\beta(P, m)\) the image of \((m, 1)\) under the canonical mapping of \(M \times P\) to \(M \otimes_A P\) ; then \(\beta(S^{-1}A, m) = 0\) . We know (Algebra, Chapter II, §6, no. 3, Corollary 4 to Proposition 7) that there exists a finitely generated submodule \(P\) of \(S^{-1}A\) containing 1 and such that \(\beta(P, m) = 0\) . For all \(t \in S\) we denote by \(A_t\) the set of \(a/t\) , where \(a \in A\) ; as \(P\) is finitely generated, there exists \(t \in S\) such that \(P \subset A_t\) (no. 1, Remark 2), whence \(\beta(A_t, m) = 0\) . The mapping \(a \mapsto a/t\) from \(A\) to \(A_t\) is surjective; let \(B\) be its kernel. It defines a surjective mapping \(h: M \otimes_A A \to M \otimes_A A_t\) , whose kernel is BM (M being identified with \(M \otimes A\) ); then \(\beta(A_t, m) = h(tm)\) and consequently \(tm\) can be expressed in the form \(\sum_{i=1}^{r} b_i m_i\) , where \(b_i \in B\) , \(m_i \in M (1 \leqslant i \leqslant r)\) . As \(b_i/t = 0\) expressed in the form \(\sum_{i=1}^{n} b_i m_i\) , where \(b_i \in \mathbf{B}\) , \(m_i \in \mathbf{M}\) ( \(1 \leqslant i \leqslant r\) ). As \(b_i / t = 0\) for \(1 \leqslant i \leqslant r\) , there exists \(t' \in \mathbf{S}\) such that \(t'b_i = 0\) for \(1 \leqslant i \leqslant r\) , whence \(t'tm = 0\) , which proves Proposition 4.

COROLLARY 1. For \(m / s = m' / s'\) in \(S^{-1}M\) , it is necessary and sufficient that there exist \(t \in S\) such that \(t(s'm - sm') = 0\) .

\[
(m / s) - (m ^ {\prime} / s ^ {\prime}) = (s ^ {\prime} m - s m ^ {\prime}) / s s ^ {\prime}.
\]

COROLLARY 2. Let M be a finitely generated A-module. For \(S^{-1}M = 0\) , it is necessary and sufficient that there exists \(s \in S\) such that \(sM = 0\) .

Without any conditions on M, clearly the relation \(sM = 0\) for some \(s \in S\) implies \(S^{-1}M = 0\) . Conversely, suppose that \(S^{-1}M = 0\) and let \((m_i)_{1 \leqslant i \leqslant n}\) be a system of generators of M; the \(m_i / 1\) generate the \(S^{-1}A\) -module \(S^{-1}M\) , hence to say that \(S^{-1}M = 0\) amounts to saying that \(m_i / 1 = 0\) for \(1 \leqslant i \leqslant n\) ; by

Proposition 4 there exist \(s_i \in S\) such that \(s_i m_i = 0\) and, taking \(s = s_1 s_2 \ldots s_n\) , \(sm_i = 0\) for all \(i\) and hence \(sM = 0\) .

COROLLARY 3. Let M be a finitely generated A-module. For an ideal a of A to be such that aM = M, it is necessary and sufficient that there exist \(a \in a\) such that \((1 + a)M = 0\) .

Clearly the relation \((1 + a)M = 0\) implies M = aM. To prove the converse we use the following lemma:

LEMMA 1. For every ideal a of A, the set S of elements \(1 + a\) , where \(a \in a\) , is a multiplicative subset of A and the set \(a'\) of elements of \(S^{-1}A\) of the form a/s, where \(a \in a\) and \(s \in S\) , is an ideal contained in the Jacobson radical of \(S^{-1}A\) .

The first assertion is obvious, as well as the fact that \(a'\) is an ideal of \(S^{-1}A\) . On the other hand, \((1/1) + (a/s) = (s + a)/s\) and \(s + a \in S\) for all \(s \in S\) and \(a \in a\) by definition of S; hence \((1/1) + (a/s)\) is invertible in \(S^{-1}A\) for all \(a/s \in a'\) , which completes the proof of the lemma (Algebra, Chapter VIII, §6, no. 3, Theorem 1).

This being so, if we set \(N = S^{-1}M\) , clearly N is a finitely generated \(S^{-1}A\) -module; if aM = M, then \(a'N = N\) and it follows that N = 0 by Nakayama's Lemma (Algebra, Chapter VIII, §6, no. 3, Corollary to Proposition 6); the corollary then follows from Corollary 2.

PROPOSITION 5. Let A, B be two rings, S a multiplicative subset of A, T a multiplicative subset of B and f a homomorphism from A to B such that \(f(S) \subset T\) . Let M be an A-module, N a B-module and u an A-linear mapping from M to N. Then there exists a unique \(S^{-1}A\) -linear mapping \(u'\) from \(S^{-1}M\) to \(T^{-1}N\) such that \(u'(m/1) = u(m)/1\) for all \(m \in M\) .

The mapping \(i_{\mathbf{N}}^{\mathbf{T}} \circ u\) from \(\mathbf{M}\) to \(\mathbf{T}^{-1}\mathbf{N}\) is A-linear. Moreover, if \(s \in \mathbf{S}\) , then \(f(s) \in \mathbf{T}\) , hence the homothety induced by \(s\) on \(\mathbf{T}^{-1}\mathbf{N}\) is bijective. The existence and uniqueness of \(u'\) then follow from Proposition 3. Then, for \(m \in \mathbf{M}\) and \(s \in \mathbf{S}\) ,

\[
u ^ {\prime} (m / s) = u (m) / f (s).\tag{6}
\]

With the same notation, let C be a third ring, U a multiplicative subset of C, g a homomorphism from B to C such that \(g(T) \subset U\) , P a C-module, v a B-linear mapping from N to P and \(v'\) the \(T^{-1}B\) -linear mapping from B to \(U^{-1}P\) associated with v. Then

\[
(v \circ u) ^ {\prime} = v ^ {\prime} \circ u ^ {\prime}\tag{7}
\]

where the left hand side is the A-linear mapping \(S^{-1}M \rightarrow U^{-1}P\) associated with \(v \circ u\) . Similarly, if \(u_{1}\) is a second A-linear mapping from M to N, then

\[
(u + u _ {1}) ^ {\prime} = u ^ {\prime} + u _ {1} ^ {\prime},\tag{8}
\]

the left-hand side being the A-linear mapping \(S^{-1}M \rightarrow T^{-1}N\) associated with \(u + u_{1}\) .

Remark (5). If, in Proposition 5, we take B = A, T = S and \(f = 1_{A}\) , it is easily seen that \(u'\) is just the mapping \(u \otimes 1: M \otimes S^{-1}A \to N \otimes S^{-1}A\) . We shall henceforth denote it by \(S^{-1}u\) ; if S is the complement of a prime ideal p of A, we write \(u_{p}\) instead of \(S^{-1}u\) .

PROPOSITION 6. Let \(f\) be a homomorphism from a ring \(\mathbf{A}\) to a ring \(\mathbf{B}\) and \(\mathbf{S}\) a multiplicative subset of \(\mathbf{A}\) . There exists a unique mapping \(j\) from \((f(\mathbf{S}))^{-1}\mathbf{B}\) to \(\mathbf{S}^{-1}\mathbf{B}\) (where \(\mathbf{B}\) is considered as an A-module by means of \(f\) ) such that \(j(b/f(s)) = b/s\) for all \(b \in \mathbf{B}\) , \(s \in \mathbf{S}\) . If \(f': \mathbf{S}^{-1}\mathbf{A} \to (f(\mathbf{S}))^{-1}\mathbf{B}\) is the ring homomorphism associated with \(f\) (no. 1, Proposition 2), then \(j \circ f' = \mathbf{S}^{-1}f\) . The mapping \(j\) is an isomorphism of the \(\mathbf{S}^{-1}\mathbf{A}\) -module structure on \((f(\mathbf{S}))^{-1}\mathbf{B}\) defined by \(f'\) onto that on \(\mathbf{S}^{-1}\mathbf{B}\) and also of the B-module structure on \((f(\mathbf{S}))^{-1}\mathbf{B}\) onto that on \(\mathbf{S}^{-1}\mathbf{B}\) (resulting from the definition \(\mathbf{S}^{-1}\mathbf{B} = (\mathbf{S}^{-1}\mathbf{A}) \otimes_{\mathbf{A}} \mathbf{B}\) ).

If \(b, b'\) are in B, \(s, s'\) in S, the conditions \(b / s = b' / s'\) and \(b / f(s) = b' / f(s')\) are equivalent, as follows from Corollary 1 to Proposition 4, which establishes the existence of \(j\) and shows that \(j\) is bijective; the uniqueness of \(j\) is obvious. Clearly \(j\) is an additive group isomorphism. If \(a \in A\) , \(b \in B\) , \(s \in S\) , \(t \in S\) , then

\[
(a / s). (b / f (t)) = f ^ {\prime} (a / s) (b / f (t)) = f (a) / f (s) (b / f (t)) = (f (a) b) / f (s t),
\]

from which its follows that \(j\) is \((\mathbf{S}^{-1}\mathbf{A})\) -linear. Clearly \(j \circ f^1 = \mathbf{S}^{-1}f\) . Finally, if \(b \in \mathbf{B}\) , \(b' \in \mathbf{B}\) , \(s \in \mathbf{S}\) , then \(j(b \cdot (b'/f(s)) = j(bb'/f(s)) = bb'/s = b \cdot (b'/s))\) , which proves the last assertion.

The mapping \(j\) of Proposition 6 is called the canonical isomorphism of \((f(\mathbf{S}))^{-1}\mathbf{B}\) onto \(\mathbf{S}^{-1}\mathbf{B}\) . These two sets are in general identified by means of \(f\) ; then \(f' = \mathbf{S}^{-1}f\) , \(i_{\mathbf{B}}^{\mathbf{S}} = i_{\mathbf{B}}^{f(\mathbf{S})}\) .

\subsection{Change of multiplicative subset}

Let A be a ring, S a multiplicative subset of A and M an A-module. If T is a multiplicative subset of A containing S, it follows from Proposition 5 of no. 2 that there exists a unique \(S^{-1}A\) -linear mapping \(i_{M}^{T,S}:S^{-1}M\to T^{-1}M\) such that \(i_{M}^{T}=i_{M}^{T,S}\circ i_{M}^{S}\) ; the mapping \(i_{M}^{T,S}\) maps the element m/s of \(S^{-1}M\) to the element m/s of \(T^{-1}M\) . It is easily verified that \(i_{M}^{T,S}=i_{A}^{T,S}\otimes l_{M}\) . If U is a third multiplicative subset of A such that \(T\subset U\) , then \(i_{M}^{U,S}=i_{M}^{U,T}\circ i_{M}^{T,S}\) ; moreover, if u: \(M\to N\) is an A-module homomorphism, the diagram

\[
\begin{array}{c} \mathrm{S} ^ {- 1} \mathrm{M} \xrightarrow {\mathrm{S} - 1 _ {u}} \mathrm{S} ^ {- 1} \mathrm{N} \\ i _ {\mathrm{M}} ^ {\mathrm{T}, \mathrm{S}} \Biggl \downarrow \qquad \qquad \qquad \qquad \qquad \qquad \qquad \Biggl \downarrow i _ {\mathrm{N}} ^ {\mathrm{T}, \mathrm{S}} \\ \mathrm{T} ^ {- 1} \mathrm{M} \xrightarrow [ \mathrm{T} - 1 _ {u} ]{} \mathrm{T} ^ {- 1} \mathrm{N} \end{array}
\]

is commutative.

PROPOSITION 7. Let A be a ring and S, T two multiplicative subsets of A. Set \(\mathbf{T}' = i(\mathbf{T}_{\mathbf{A}}^{\mathbf{s}})\) .

\begin{enumerate}
\def\labelenumi{(\roman{enumi})}
\tightlist
\item
  There exists a unique isomorphism \(j\) from the ring (ST) \(^{-1}\) A onto the ring \(\mathbf{T}^{\prime -1}(\mathbf{S}^{-1}\mathbf{A})\) such that the diagram

\[
\begin{array}{c} \mathrm{A} \xrightarrow {i _ {\mathrm{A}} ^ {\mathrm{S}}} \mathrm{S} ^ {- 1} \mathrm{A} \\ i _ {\mathrm{M}} ^ {\mathrm{ST}} \Big \downarrow \qquad \qquad \qquad \qquad \qquad \qquad \qquad \Big \downarrow i _ {\mathrm{S} - 1 _ {\mathrm{A}}} ^ {\mathrm{T} ^ {\prime}} \\ (\mathrm{ST}) ^ {- 1} \mathrm{A} \xrightarrow [ j ]{} \mathrm{T} ^ {\prime - 1} (\mathrm{S} ^ {- 1} \mathrm{A}) \end{array}
\]

is commutative.

\item
  Let M be an A-module. There exists an \((\mathrm{ST})^{-1}\mathrm{A}\) -isomorphism k from the \((\mathrm{ST})^{-1}\mathrm{A}\) -module \((\mathrm{ST})^{-1}\mathrm{M}\) onto the \(\mathrm{T}^{\prime-1}(\mathrm{S}^{-1}\mathrm{A})\) -module \(\mathrm{T}^{\prime-1}(\mathrm{S}^{-1}\mathrm{M})\) such that the diagram
\end{enumerate}

\[
\begin{array}{c} \mathrm{M} \qquad \xrightarrow {i _ {\mathrm{M}} ^ {\mathrm{S}}} \mathrm{S} ^ {- 1} \mathrm{M} \\ i _ {\mathrm{A}} ^ {\mathrm{ST}} \Big \downarrow \qquad \qquad \qquad \qquad \qquad \Big \downarrow i _ {\mathrm{S}} ^ {\mathrm{T} ^ {\prime} - 1 _ {\mathrm{M}}} \\ (\mathrm{ST}) ^ {- 1} \mathrm{M} \xrightarrow [ k ]{} \mathrm{T} ^ {\prime - 1} (\mathrm{S} ^ {- 1} \mathrm{M}) \end{array}
\]

is commutative.

\begin{enumerate}
\def\labelenumi{(\roman{enumi})}
\tightlist
\item
  We use the definition of \((\mathrm{ST})^{-1}\mathrm{A}\) as the solution of a universal mapping problem. Let B be a ring and f a homomorphism from A to B such that \(f(\mathrm{ST})\) consists of invertible elements. As \(f(\mathrm{S})\) consequently consists of invertible elements, there exists a unique homomorphism \(f': S^{-1}A \to B\) such that \(f = f' \circ i_{A}^{S}\) (no. 1, Proposition 1). For all \(t \in T\) , \(f'(i_{\mathbf{A}}^{\mathrm{S}}(t)) = f(t)\) is invertible in B by hypothesis, hence \(f'(\mathrm{T}')\) consists of invertible elements; then there exists, by no. 1, Proposition 1, a unique homomorphism \(f''\) from \(T'^{-1}(S^{-1}A)\) to B such that \(f' = f'' \circ i_{S^{-1}A}^{T'}\) , whence \(f = f'' \circ u\) , setting \(u = i_{S^{-1}A}^{T'} \circ i_{A}^{S}\) .

Moreover, if \(f_1'' \colon \mathrm{T}^{-1}(\mathrm{S}^{-1}\mathrm{A}) \to \mathrm{B}\) is a second homomorphism such that \(f_1'' \circ u = f\) , then \((f_1'' \circ i_{\mathrm{S}^{-1}\mathrm{A}}^{\mathrm{T}'} \circ i_{\mathrm{A}}^{\mathrm{S}} = (f'' \circ i_{\mathrm{S}^{-1}\mathrm{A}}^{\mathrm{T}'} \circ i_{\mathrm{A}}^{\mathrm{S}} \text{, whence } f_1'' \circ i_{\mathrm{S}^{-1}\mathrm{A}}^{\mathrm{T}'} = f'' \circ i_{\mathrm{S}^{-1}\mathrm{A}}^{\mathrm{T}'} \text{ and consequently } f_1'' = f''\) .

As the images under \(u\) of the elements of ST in \(\mathbf{T}^{-1}(\mathbf{S}^{-1}\mathbf{A})\) are invertible, the ordered pair \((\mathbf{T}^{-1}(\mathbf{S}^{-1}\mathbf{A}), u)\) is a solution of the universal mapping problem (relative to A and ST) considered in no. 1. This shows the existence and uniqueness of \(j\) .

\item
  The proof is completely analogous with that of (i), using in this case no. 2, Proposition 3, and is left to the reader.
\end{enumerate}

PROPOSITION 8. Let A be a ring and S, T two multiplicative subsets of A such that \(\mathbf{S} \subset \mathbf{T}\) . The following properties are equivalent:

\begin{enumerate}
\def\labelenumi{(\alph{enumi})}
\item
  The homomorphism \(i_{\mathbf{A}}^{\mathrm{T},\mathrm{S}}: \mathrm{S}^{-1}\mathrm{A} \to \mathrm{T}^{-1}\mathrm{A}\) is bijective.
\item
  For every A-module M, the homomorphism \(i_{\mathbf{A}}^{\mathrm{T},\mathrm{S}}: \mathrm{S}^{-1}\mathrm{M} \to \mathrm{T}^{-1}\mathrm{M}\) is bijective.
\item
  For all \(t \in \mathbf{T}\) , there exists \(a \in \mathbf{A}\) such that \(at \in \mathbf{S}\) (in other words, every element of \(\mathbf{T}\) divides an element of \(\mathbf{S}\) ).
\item
  Every prime ideal which meets T meets S.
\end{enumerate}

It has been seen above that \(i_{\mathbf{A}}^{\mathrm{T},\mathrm{S}} = 1_{\mathbb{M}} \otimes i_{\mathbf{A}}^{\mathrm{T},\mathrm{S}}\) , which immediately proves the equivalence of (a) and (b). Set \(\mathrm{T}' = i_{\mathbf{A}}^{\mathrm{S}}(\mathrm{T})\) ; then (Proposition 7) \(\mathrm{T}^{-1}\mathbf{A}\) is identified with \(\mathrm{T}'^{-1}(\mathrm{S}^{-1}\mathbf{A})\) and (a) is equivalent to saying that the elements of \(\mathrm{T}'\) are invertible in \(\mathrm{S}^{-1}\mathbf{A}\) (no. 1, Remark 5). Now, to say that \((t/1)(a/s) = 1/1\) ( \(t \in \mathrm{T}\) , \(a \in \mathrm{A}\) , \(s \in \mathrm{S}\) ) means that there exists \(s' \in \mathrm{S}\) such that \(tas' = ss'\) , which shows the equivalence of (a) and (c). We show that (d) implies (c). Let \(t\) be an element of \(\mathrm{T}\) and suppose that \(t/1\) is not invertible in \(\mathrm{S}^{-1}\mathbf{A}\) ; then there exists a maximal ideal \(m'\) of \(\mathrm{S}^{-1}\mathbf{A}\) containing \(t/1\) (Algebra, Chapter I, § 8, no. 7, Theorem 2) and \(p = (i_{\mathbf{A}}^{\mathrm{S}})^{-1}(m')\) is a prime ideal of \(\mathbf{A}\) containing \(t\) and not meeting \(\mathbf{S}\) (since the image under \(i_{\mathbf{A}}^{\mathrm{S}}\) of an element of \(\mathbf{S}\) is invertible). Conversely, if there exists a prime ideal \(p\) which meets \(\mathrm{T}\) without meeting \(\mathrm{S}\) , then no element of \(p \cap \mathrm{T}\) can divide an element of \(\mathrm{S}\) ; this proves that (c) implies (d) and completes the proof.

It follows from Proposition 8 that, amongst the multiplicative subsets T of A containing S and satisfying the equivalent conditions of Proposition 8, there exists a greatest, consisting of all the elements of A which divide an element of S (cf.~Exercise 1).

PROPOSITION 9. Let I be a right directed preordered set, \((\mathrm{S}_{\alpha})_{\alpha \in \mathbf{I}}\) an increasing family of multiplicative subsets of a ring A and \(S = \bigcup_{\alpha \in I} S_{\alpha}\) . We write \(\rho_{\beta \alpha} = i_{A}^{S_{\beta}, S_{\alpha}}\) for \(\alpha \leqslant \beta\) , \(\rho_{\alpha} = i_{A}^{S, S_{\alpha}}\) . Then \((S_{\alpha}^{-1}A, \rho_{\beta \alpha})\) is a direct system of rings and, if, for all \(\alpha \in I\) , \(\rho_{\alpha}'\) is the canonical mapping of \(\mathbf{S}_{\alpha}^{-1}\mathbf{A}\) to \(\lim_{\longrightarrow}\mathbf{S}_{\alpha}^{-1}\mathbf{A}\) , there exists a unique isomorphism \(j\) from \(\lim_{\longrightarrow}\mathbf{S}_{\alpha}^{-1}\mathbf{A}\) to \(\mathbf{S}^{-1}\mathbf{A}\) such that \(j\circ \rho_{\alpha}^{\prime} = \rho_{\alpha}\) for all \(\alpha \in \mathbf{I}\) .

For \(\alpha \leqslant \beta \leqslant \gamma\) , \(\rho_{\gamma \alpha} = \rho_{\gamma \beta} \circ \rho_{\beta \alpha}\) (no. 1, Corollary 3 to Proposition 2), hence \((S_{\alpha}^{-1}A, \rho_{\beta \alpha})\) is a direct system. We write \(A' = \varinjlim S_{\alpha}^{-1}A\) ; as \(\rho_{\alpha} = \rho_{\beta} \circ \rho_{\beta \alpha}\) for \(\alpha \leqslant \beta\) (no. 1, Corollary 3 to Proposition 2), \((\rho_{\alpha})\) is a direct system of homomorphisms and \(j = \varinjlim \rho_{\alpha}\) is the unique homomorphism from \(A'\) to \(S^{-1}A\) such that \(j \circ \rho_{\alpha}' = \rho_{\alpha}\) for all \(\alpha \in I\) . The homomorphisms \(\rho_{\alpha}' \circ i_{A^{\alpha}}^{S}: A \to A'\) are all equal, for \(\rho_{\beta \alpha} \circ i_{A^{\alpha}}^{S} = i_{A^{\beta}}^{S}\) for \(\alpha \leqslant \beta\) ; let \(u\) be their common value. Clearly the elements of \(u(S)\) are invertible in \(A'\) , which shows that there exists a homomorphism \(h: S^{-1}A \to A'\) such that \(h \circ i_{A}^{S} = u\) (no. 1, Proposition 1). Then

\[
j \circ h \circ i _ {\mathrm{S}} ^ {\mathrm{A}} = j \circ u = j \circ \rho_ {\alpha} ^ {\prime} \circ i _ {\mathrm{S} ^ {\alpha}} ^ {\mathrm{A}} = \rho_ {\alpha} \circ i _ {\mathrm{A} ^ {\alpha}} ^ {\mathrm{S}} = i _ {\mathrm{A}} ^ {\mathrm{S}}
\]

for all \(\alpha \in \mathbf{I}\) and consequently \(j \circ h\) is the identity automorphism of \(\mathbf{S}^{-1}\mathbf{A}\) . On the other hand, for all \(\alpha \in \mathbf{I}\) ,

\[
h \circ j \circ \rho_ {\alpha} ^ {\prime} \circ i _ {A ^ {\alpha}} ^ {S} = h \circ \rho_ {\alpha} \circ i _ {A ^ {\alpha}} ^ {S} = h \circ i _ {A} ^ {S} = u = \rho_ {\alpha} ^ {\prime} \circ i _ {A ^ {\alpha}} ^ {S},
\]

whence \(h \circ j \circ \rho_{\alpha}' = \rho_{\alpha}'\) for all \(\alpha \in I\) ; it follows that \(h \circ j\) is the identity automorphism of \(A'\) and consequently \(j\) is an isomorphism.

COROLLARY. Under the hypotheses of Proposition 9, let M be an A-module. We write \(f_{\beta \alpha} = i_{\mathrm{M}}^{\mathrm{S}_{\beta},\mathrm{S}_{\alpha}}\) for \(\alpha \leqslant \beta, f_{\alpha} = i_{\mathrm{M}}^{\mathrm{S},\mathrm{S}_{\alpha}}\) for all \(\alpha \in \mathbf{I}\) and let \(f_{\alpha}'\) be the canonical mapping from \(\mathrm{S}_{\alpha}^{-1}\mathrm{M}\) to \(\lim \mathrm{S}_{\alpha}^{-1}\mathrm{M}\) ; then there exists an \(\mathrm{S}^{-1}\mathrm{A}\) -isomorphism \(g\) of \(\mathrm{S}^{-1}\mathrm{M}\) onto \(\lim \mathrm{S}_{\alpha}^{-1}\mathrm{M}\) such that \(g \circ f_{\alpha} = f_{\alpha}'\) for all \(\alpha \in \mathbf{I}\) .

The corollary follows immediately from the definitions \(S_{\alpha}^{-1}M = M \otimes_{A} S_{\alpha}^{-1}A\) and \(S^{-1}M = M \otimes_{A} S^{-1}A\) and the fact that taking direct limits commutes with tensor products (Algebra, Chapter II, § 6, no. 3, Proposition 7).

\subsection{Properties of modules of fractions}

Throughout this no., A denotes a ring and S a multiplicative subset of A.

Let \((\mathbf{M}_{\alpha}, \phi_{\beta\alpha})\) be a direct system of A-modules; then \((\mathrm{S}^{-1}\mathbf{M}_{\alpha}, \mathrm{S}^{-1}\phi_{\beta\alpha})\) is a direct system of \(S^{-1}A\) -modules and the fact that taking direct limits commutes with tensor products (Algebra, Chapter II, § 6, no. 3, Proposition 7) allows us to define a canonical isomorphism

\[
\lim _ {\longrightarrow} \left(\mathrm{S} ^ {- 1} \mathrm{M} _ {\alpha}\right)\rightarrow \mathrm{S} ^ {- 1} \lim _ {\longrightarrow} \mathrm{M} _ {\alpha}.
\]

Similarly, the fact that taking direct sums commutes with tensor products (Algebra, Chapter II, § 3, no. 7, Proposition 7) allows us to define for every family \((\mathbf{M}_t)_{i\in I}\) of A-modules a canonical isomorphism

\[
\bigoplus_ {i \in I} S ^ {- 1} M _ {i} \rightarrow S ^ {- 1} \bigoplus_ {i \in I} M _ {i}.
\]

Finally we note that, if an A-module M is the sum of a family \((\mathbf{N}_{\iota})_{\iota\in I}\) of submodules, \(S^{-1}M\) is the sum of the family of sub- \(S^{-1}A\) -modules generated by the \(i_{M}^{S}(N_{\iota})\) . Then it follows that, if M is a finitely generated A-module (resp. finitely presented A-module), \(S^{-1}M = S^{-1}A \otimes_{A} M\) is a finitely generated \(S^{-1}A\) -module (resp. finitely presented \(S^{-1}A\) -module).

THEOREM 1. The ring \(\mathbf{S}^{-1}\mathbf{A}\) is a flat A-module (Chapter I, § 2, no. 1, Definition 2).

If \(u: \mathbf{M}' \to \mathbf{M}\) is an injective homomorphism of A-modules, it is necessary to establish that \(\mathbf{S}^{-1}u: \mathbf{S}^{-1}\mathbf{M}' \to \mathbf{S}^{-1}\mathbf{M}\) is injective. Now, if \(m'/s\) ( \(m' \in \mathbf{M}'\) , \(s \in \mathbf{S}\) ) is such that \(u(m')/s = 0\) , this implies the existence of an \(s' \in \mathbf{S}\) such that \(s'u(m') = 0\) (no. 2, Proposition 4) or \(u(s'm') = 0\) ; as \(u\) is injective, it follows that \(s'm' = 0\) , whence \(m'/s = 0\) .

The fact that \(\mathbf{S}^{-1}\mathbf{A}\) is a flat A-module allows us to apply to it the results of Chapter I, § 2. In particular:

\begin{enumerate}
\def\labelenumi{(\arabic{enumi})}
\tightlist
\item
  If M is an A-module and N a submodule of M, \(S^{-1}N\) is canonically identified with a submodule of \(S^{-1}M\) generated by \(i_{\mathrm{M}}^{\mathrm{S}}(\mathrm{N})\) (Chapter I, § 2, no. 3, Remark 2); with this identification, \(S^{-1}(\mathrm{M}/\mathrm{N})\) is identified with \((\mathrm{S}^{-1}\mathrm{M})/(\mathrm{S}^{-1}\mathrm{N})\) and, if P is a second submodule of M, then
\end{enumerate}

\[
\mathrm{S} ^ {- 1} (\mathrm{N} + \mathrm{P}) = \mathrm{S} ^ {- 1} \mathrm{N} + \mathrm{S} ^ {- 1} \mathrm{P}, \quad \mathrm{S} ^ {- 1} (\mathrm{N} \cap \mathrm{P}) = \mathrm{S} ^ {- 1} \mathrm{N} \cap \mathrm{S} ^ {- 1} \mathrm{P}
\]

(Chapter I, § 2, no. 6, Proposition 2).

\begin{enumerate}
\def\labelenumi{(\arabic{enumi})}
\setcounter{enumi}{1}
\tightlist
\item
  If M is a finitely generated A-module, then
\end{enumerate}

\[
\mathbf {S} ^ {- 1} \operatorname{Ann} (\mathbf {M}) = \operatorname{Ann} (\mathbf {S} ^ {- 1} \mathbf {M})\tag{9}
\]

(Chapter I, § 2, no. 10, Corollary 2 to Proposition 12).

PROPOSITION 10. Let M be an A-module. For every submodule N' of the S \(^{-1}\) A-module S \(^{-1}\) M, let \(\phi(N')\) be the inverse image of N' under \(i_{M}^{8}\) . Then:

\begin{enumerate}
\def\labelenumi{(\roman{enumi})}
\item
  \(\mathbf{S}\phi(\mathbf{N}^{\prime}) = \mathbf{N}^{\prime}.\)
\item
  For every submodule \(\mathbf{N}\) of \(\mathbf{M}\) , the submodule \(\phi (\mathrm{S}^{-1}\mathrm{N})\) of \(\mathbf{M}\) consists of those \(m\in \mathbf{M}\) for which there exists \(s\in \mathbf{S}\) such that \(sm\in \mathbf{N}\) .
\item
  \(\phi\) is an isomorphism (for the orderings defined by inclusion) of the set of sub- \(\mathbf{S}^{-1}\mathbf{A}\) -modules of \(\mathbf{S}^{-1}\mathbf{M}\) onto the set of submodules \(\mathbf{Q}\) of \(\mathbf{M}\) which satisfy the following condition:

(MS) If \(sm \in \mathbb{Q}\) , where \(s \in \mathbb{S}\) , \(m \in \mathbb{M}\) , then \(m \in \mathbb{Q}\) .

\end{enumerate}
Obviously \(\mathrm{S}^{-1}\phi(\mathrm{N}')\subset\mathrm{N}'\) ; conversely, if \(n' = m/s \in N'\) , then \(m/1 \in N'\) , hence \(m \in \phi(\mathrm{N}')\) and consequently \(n' \in \mathrm{S}^{-1}(\phi(\mathrm{N'}))\) ; whence (i). For an element \(m \in M\) to be such that \(m \in \phi(\mathrm{S}^{-1}\mathrm{N})\) , it is necessary and sufficient that \(m/1 \in S^{-1}N\) , that is that there exist \(s \in S\) and \(n \in N\) such that \(m/1 = n/s\) ; this means that there exists \(s' \in S\) such that \(s'sm = s'n \in N\) , whence (ii). Finally, the relation \(sm \in \phi(\mathrm{N}')\) is equivalent by definition to \(sm/1 \in N'\) and as s/1 is invertible in \(S^{-1}A\) , this implies \(m/1 \in N'\) , or \(m \in \phi(\mathrm{N}')\) , hence \(\phi(\mathrm{N}')\) satisfies condition (MS); on the other hand, it follows from (ii) that, if N satisfies (MS), then \(\phi(\mathrm{S}^{-1}\mathrm{N}) = N\) , which completes the proof of (iii).

The submodule \(\phi(\mathrm{S}^{-1}\mathrm{N})\) is called the saturation of N in M with respect to S, and the submodules satisfying condition (MS) (and hence equal to their saturations) are said to be saturated with respect to S. The submodule \(\phi(\mathrm{S}^{-1}\mathrm{N})\) is the kernel of the composite homomorphism

\[
\mathrm{M} \xrightarrow {h} \mathrm{M} / \mathrm{N} \xrightarrow {i _ {\mathrm{M} / \mathrm{N}} ^ {\mathrm{S}}} \mathrm{S} ^ {- 1} \mathrm{M} / \mathrm{S} ^ {- 1} \mathrm{N}
\]

where h is the canonical homomorphism, as follows from the commutativity of the diagram

\[
\begin{array}{c c c} \mathrm{M} & \stackrel {{h}} {{\longrightarrow}} & \mathrm{M/N} \\ i _ {\mathrm{M}} ^ {\mathrm{S}} \Big \downarrow & & \Big \downarrow i _ {\mathrm{M/N}} ^ {\mathrm{S}} \\ \mathrm{S} ^ {- 1} \mathrm{M} & \xrightarrow [ \mathrm{S} ^ {- 1 h} ]{} \mathrm{S} ^ {- 1} \mathrm{M/S} ^ {- 1} \mathrm{N} \end{array}
\]

If S is the complement in A of a prime ideal p, \(\phi(S^{-1}N)\) is also called the saturation of N in M with respect to p.

COROLLARY 1. Let \(N_{1}\) , \(N_{2}\) be two submodules of an A-module M. For \(S^{-1}N_{1} \subset S^{-1}N_{2}\) , it is necessary and sufficient that the saturation of \(N_{1}\) with respect to S be contained in that of \(N_{2}\) .

COROLLARY 2. If M is a Noetherian (resp. Artinian) A-module, S \(^{-1}\) M is a Noetherian (resp. Artinian) S \(^{-1}\) A-module. In particular, if the ring A is Noetherian (resp. Artinian), so is the ring S \(^{-1}\) A.

\subsection{Ideals in a ring of fractions}

PROPOSITION 11. Let A be a ring and S a multiplicative subset of A. For every ideal \(\mathfrak{b}'\) of \(\mathbf{S}^{-1}\mathbf{A}\) , let \(\mathfrak{b} = (i_{\mathbf{A}}^{\mathbf{S}})^{-1}(\mathfrak{b}')\) be such that \(\mathfrak{b}' = \mathbf{S}^{-1}\mathfrak{b}\) .

\begin{enumerate}
\def\labelenumi{(\roman{enumi})}
\item
  Let \(f\) be the canonical homomorphism \(\mathbf{A} \to \mathbf{A} / \mathfrak{b}\) . The homomorphism from \(\mathbf{S}^{-1}\mathbf{A}\) to \((f(\mathbf{S}))^{-1}(\mathbf{A} / \mathfrak{b})\) canonically associated with \(f\) (no. 1, Proposition 2) is surjective and its kernel is \(\mathfrak{b}'\) , which defines, by taking quotients, a canonical isomorphism of \((\mathbf{S}^{-1}\mathbf{A}) / \mathfrak{b}'\) onto \((f(\mathbf{S}))^{-1}(\mathbf{A} / \mathfrak{b})\) . Moreover, the canonical homomorphism from \(\mathbf{A} / \mathfrak{b}\) to \((f(\mathbf{S}))^{-1}(\mathbf{A} / \mathfrak{b})\) is injective.
\item
  The mapping \(\mathfrak{b}' \mapsto \mathfrak{b} = (i_{\mathbf{M}}^{\mathbf{S}})^{-1}(\mathfrak{b}')\) , restricted to the set of maximal (resp. prime) ideals of \(\mathbf{S}^{-1}\mathbf{A}\) , is an isomorphism (with respect to inclusion) of this set onto the set of ideals of \(\mathbf{A}\) which are maximal among those which do not meet \(\mathbf{S}\) (resp. the set of prime ideals of \(\mathbf{A}\) not meeting \(\mathbf{S}\) ).
\item
  If \(\mathfrak{q}'\) is a prime ideal of \(S^{-1}A\) and \(\mathfrak{q} = (i_A^S)^{-1}(\mathfrak{q}')\) , there exists an isomorphism of the ring of fractions \(A_{\mathfrak{q}}\) onto the ring \((S^{-1}A)_{\mathfrak{q}}\) , which maps \(a/b\) to \((a/1)/(b/1)\) , where \(a \in A\) , \(b \in A - \mathfrak{q}\) .

\end{enumerate}

\begin{enumerate}
\def\labelenumi{(\roman{enumi})}
\item
  \((f(S))^{-1}(A/b)\) can be identified with \(S^{-1}(A/b)\) by means of the canonical isomorphism between these two modules (no. 2, Proposition 6). The exact sequence \(0 \rightarrow b \rightarrow A \rightarrow A/b \rightarrow 0\) then induces an exact sequence

\[
0 \rightarrow \mathrm{S} ^ {- 1} \mathfrak {b} \rightarrow \mathrm{S} ^ {- 1} \mathrm{A} \rightarrow \mathrm{S} ^ {- 1} (\mathrm{A} / \mathfrak {b}) \rightarrow 0
\]

(no. 4, Theorem 1) whose existence proves the first assertion of (i), taking account of the fact that \(\mathfrak{b}' = \mathrm{S}^{-1}\mathfrak{b}\) . Since \(\mathfrak{b}\) is saturated with respect S, the conditions \(a \in \mathbf{A}\) , \(s \in \mathbf{S}\) , \(as \in \mathfrak{b}\) imply \(a \in \mathfrak{b}\) ; the homothety of ratio \(s\) on \(\mathbf{A}/\mathfrak{b}\) is then injective, which proves the second assertion of (i).

\item
  We note first that the relation \(b' = S^{-1}A\) is equivalent to the relation \(b \cap S \neq \varnothing\) , the latter expressing the fact that \(b'\) contains invertible elements of \(S^{-1}A\) . It follows from no. 4, Proposition 10 (iii) that \(b' \mapsto b = (i_{\mathrm{A}}^{\mathrm{S}})^{-1}(b')\) is an isomorphism (with respect to inclusion) of the set of ideals of \(S^{-1}A\) distinct from \(S^{-1}A\) onto the set F of ideals of A not meeting S and satisfying condition (MS) of Proposition 10. If \(b'\) is maximal (resp. prime), clearly \(b'\) is maximal in F (resp. prime) and conversely (by (i)). On the other hand, if r is an ideal of A disjoint from S, its saturation \(r_{1}\) with respect to S is an ideal of A containing r and disjoint from S: for no element \(a \in S\) can satisfy \(sa \in r\) for some \(s \in S\) , since it would follow that \(sa \in r \cap S\) . We conclude that, if r is maximal among the ideals of A meeting S, it is maximal in F. Similarly, if r is a prime ideal not meeting S, it satisfies condition (MS) of no. 4, Proposition 10 by definition of prime ideals and hence belongs to F This completes the proof of (ii).
\item
  Suppose that \(q'\) is prime and such that q is also prime. The set \(T = A - q\) is a multiplicative subset of A which contains S, whence \(ST = T\) . We write \(T' = i_{\mathbf{A}}^{\mathbf{S}}(\mathbf{T})\) ; it follows from no. 3, Proposition 7 (i) that there exists a unique isomorphism j of \(T^{-1}A = A_{q}\) onto \(T'^{-1}(S^{-1}A)\) such that

\[
j (a / b) = (a / 1) / (b / 1),
\]

where \(a \in \mathbf{A}\) and \(b \in \mathbf{T}\) . On the other hand \(\mathbf{T}'\) obviously does not meet \(\mathfrak{q}'\) ; conversely, let \(a / s \in \mathrm{S}^{-1}\mathrm{A}\) ; since \(1 / s\) is invertible in \(\mathrm{S}^{-1}\mathrm{A}\) , the condition \(a / s \notin \mathfrak{q}'\) is equivalent to \(i_{\mathbf{A}}^{\mathbf{S}}(a) = a / 1 \notin \mathfrak{q}'\) and hence to \(a \notin \mathfrak{q}\) ; it follows that \(\mathrm{S}^{-1}\mathrm{A} - \mathfrak{q}' = \mathrm{S}^{-1}\mathrm{T}'\) and hence, by Proposition 8 of no. 3, \(\mathrm{T}'^{-1}(\mathrm{S}^{-1}\mathrm{A}) = (\mathrm{S}^{-1}\mathrm{A})_{\mathfrak{q}'}\) .

The isomorphism defined in (iii) is called canonical. If A is an integral domain, the canonical isomorphisms of \(A_{q}\) and \((\mathrm{S}^{-1}\mathrm{A})_{q'}\) onto subrings of the field of fractions K of A have the same image.
\end{enumerate}

Remark. For an ideal a of A to satisfy \(S^{-1}a = S^{-1}A\) (or, what amounts to the same thing by no. 4, Theorem 1, \(S^{-1}(A/a) = 0\) ), it is necessary and sufficient that \(a \cap S \neq \varnothing\) , as follows immediately from the definitions.

COROLLARY 1. Let A be a ring and S a multiplicative subset of A. Every ideal p of A which is maximal among those which do not meet S is prime.

By Proposition 11, the hypothesis on p means that \(\mathfrak{p} = (i_{\mathrm{A}}^{\mathrm{S}})^{-1}(\mathfrak{m}')\) , where \(m'\) is a maximal ideal of \(S^{-1}A\) ; as \(m'\) is prime, so is p.

COROLLARY 2. Let A be a ring and S a multiplicative subset of A. For every ideal a of A not meeting S there exists a prime ideal containing a and not meeting S.

\(S^{-1}a \neq S^{-1}A\) (Remark) and hence there exists a maximal ideal of \(S^{-1}A\) containing \(S^{-1}a\) (Algebra, Chapter I, § 8, no. 7, Theorem 2) and the corollary follows from Proposition 11 (ii).

COROLLARY 3. Let A, B be two rings, \(\rho\) a homomorphism from A to B and \(\mathfrak{p}\) a prime ideal of A. For there to exist a prime ideal \(\mathfrak{p}'\) of B such that \(\overline{\rho}^{1}(\mathfrak{p}') = \mathfrak{p}\) , it is necessary and sufficient that \(\overline{\rho}^{1}(\mathrm{B}\rho(\mathfrak{p})) = \mathfrak{p}\) .

If there exists an ideal \(p'\) of \(B\) such that \(\overline{\rho}^{1}(p') = p\) , then \(\rho(p) \subset p'\) , whence \(B\rho(p) \subset p'\) and \(\overline{\rho}^{1}(B\rho(p)) \subset \overline{\rho}^{1}(p') \subset p\) ; as the converse inclusion is obvious, \(\overline{\sigma}^{1}(B\rho(p)) = p\) . Conversely, suppose that \(\overline{\rho}^{1}(B\rho(p)) = p\) and consider the multiplicative subset \(S = \rho(A - p)\) of \(B\) ; the hypothesis shows that

\[
\mathrm{S} \cap \mathrm{B} _ {\rho} (\mathfrak {p}) = \varnothing ;
\]

by Corollary 2 there exists a prime ideal \(\mathfrak{p}'\) of B containing \(\mathrm{B}\varphi(\mathfrak{p})\) and not meeting S; then \(\overline{\rho}^{1}(\mathfrak{p}')\) contains \(\mathfrak{p}\) and cannot contain any element of A - \(\mathfrak{p}\) and hence is equal to \(\mathfrak{p}\) .

COROLLARY 4. Let A, B be two rings and \(\rho\) a homomorphism from A to B.

\begin{enumerate}
\def\labelenumi{(\roman{enumi})}
\item
  Suppose that there exists a B-module E such that \(\rho_{*}(\mathbf{E})\) is a faithfully flat A-module. Then, for every prime ideal \(\mathfrak{p}\) of A, there exists a prime ideal \(\mathfrak{p}'\) of B such that \(\overline{\rho}^{-1}(\mathfrak{p}') = \mathfrak{p}\) .
\item
  Conversely, suppose that B is a flat A-module. Then, if, for every prime ideal p of A, there exists an ideal \(p'\) of B such that \(\overline{\rho}^{-1}(p') = p\) , B is a faithfully flat A-module.
\end{enumerate}

\begin{enumerate}
\def\labelenumi{(\roman{enumi})}
\item
  The hypothesis implies that, for every ideal \(\mathfrak{a}\) of \(\mathbf{A}\) , \(\bar{\rho}^{1}(\mathrm{B}\varphi (\mathfrak{a})) = \mathfrak{a}\) (Chapter I, § 3, no. 5, Proposition 8 (ii)) and it is sufficient to apply Corollary 3.

\item
  It is sufficient to show that, for every maximal ideal m of A, there exists a maximal ideal \(m'\) of B such that \(\bar{\rho}^{-1}(m') = m\) (Chapter I, §2, no. 5, Proposition 9 (e)). Now there exists by hypothesis an ideal q of B such that \(\bar{\rho}^{-1}(q) = m\) ; as \(q \neq B\) , there exists a maximal ideal \(m'\) of B containing q and consequently \(\bar{\rho}^{-1}(m') \supset m\) ; but, as \(\bar{\rho}^{-1}(m')\) cannot contain 1, \(\bar{\rho}^{-1}(m') = m\) .
\end{enumerate}

COROLLARY 5. Let A be a ring, S a multiplicative subset of A and B a ring such that \(i_{\mathbf{A}}^{\mathbf{S}}(\mathbf{A}) \subset \mathbf{B} \subset \mathbf{S}^{-1}\mathbf{A}\) . Let \(\mathfrak{q}\) be a prime ideal of B such that the prime ideal \(\mathfrak{p} = (i_{\mathbf{A}}^{\mathbf{S}})^{-1}(\mathfrak{q})\) of A does not meet S and let \(\mathfrak{p}'\) be the prime ideal \(\mathbf{S}^{-1}\mathfrak{p}\) of \(\mathbf{S}^{-1}\mathbf{A}\) . Then \(\mathfrak{p}' \cap \mathbf{B} = \mathfrak{q}\) .

Let \(S' = i_{\mathbf{A}}^{\mathbf{S}}(\mathbf{S})\) ; a canonical isomorphism has been defined from \(S'^{-1}\mathbf{B}\) to \(S^{-1}\mathbf{A}\) (no. 1, Corollary 4 to Proposition 2); we identify these two rings by means of this isomorphism. As \(q \cap S' = \varnothing\) , \(q' = S'^{-1}q\) is the unique prime ideal of \(S^{-1}\mathbf{A} = S'^{-1}\mathbf{B}\) such that \(q' \cap B = (i_{\mathbf{A}}^{\mathbf{S}})^{-1}(q') = q\) (Proposition 11 (ii)), whence \((i_{\mathbf{A}}^{\mathbf{S}})^{-1}(q') = p\) ; consequently \(q' = p'\) (Proposition 11 (ii)).

In the notation of Corollary 5, there are canonical isomorphisms of \(\mathbf{A}_{\mathfrak{p}}\) and \(\mathbf{B}_{\mathfrak{q}}\) onto \((\mathrm{S}^{-1}\mathrm{A})_{\mathfrak{p}}\) . (Proposition 11 (iii)) and hence a canonical isomorphism \(\mathbf{A}_{\mathfrak{p}} \to \mathbf{B}_{\mathfrak{q}}\) .

\subsection{Nilradical and minimal prime ideals}

In a (commutative) ring A the set of nilpotent elements is an ideal, for if x, y are elements of A such that \(x^{m} = y^{n} = 0\) , then \((x + y)^{m+n} = 0\) by the binomial theorem.

DEFINITION 4. The ideal of nilpotent elements of a (commutative) ring A is called the nilradical of A. If a is an ideal of A, the inverse image, under the canonical mapping \(A \rightarrow A/a\) , of the nilradical of A/a is called the radical of a.

We often denote by \(\mathbf{r}(\mathfrak{a})\) the radical of an ideal a of A.

To say that an element \(x \in A\) belongs to the radical of \(a\) means therefore that there exists an integer \(n > 0\) such that \(x^n \in a\) . If \(f\) is a homomorphism from \(A\) to a ring \(B\) and \(b\) is an ideal of \(B\) , the radical of \(f(b)\) is the inverse image under \(f\) of the radical of \(b\) , for to say that \(x^n \in f(b)\) means that \((f(x))^n \in b\) .

The nilradical of a ring A is contained in its Jacobson radical (Algebra, Chapter VIII, § 6, no. 3, Corollary 3 to Theorem 1) but may be distinct from it; it is always equal to it if A is Artinian (Algebra, Chapter VIII, § 6, no. 4, Theorem 3).

We say that a prime ideal p of a ring A is a minimal prime ideal if it is minimal in the set of prime ideals of A ordered by inclusion.

PROPOSITION 12. Let \(\mathfrak{p}\) be a minimal prime ideal of a ring \(\mathbf{A}\) . For all \(x \in \mathfrak{p}\) , there exists \(s \in \mathbf{A} - \mathfrak{p}\) and an integer \(k > 0\) such that \(sx^k = 0\) .

The set S of elements of the form \(sx^{k}\) (k an integer \textgreater0, \(s \in A - p\) ) is a multiplicative subset of A. If \(0 \notin S\) , there would exist a prime ideal \(p'\) not meeting S (no. 5, Corollary 2 to Proposition 11). Then \(p' \subset p\) and \(p' \neq p\) since \(x \notin p'\) , contrary to the hypothesis that p is minimal.

PROPOSITION 13. The nilradical of a ring A is the intersection of all the prime ideals of A and it is also the intersection of the minimal prime ideals of A.

Clearly, if \(x \in A\) is nilpotent, it is contained in every prime ideal of A ( \(\S 1\) , no. 1, Definition 1). Conversely, let x be a non-nilpotent element of A; the set S of \(x^{k}\) (k an integer \(\geqslant 0\) ) is then a multiplicative subset of A not containing 0 and hence there exists a prime ideal p of A not meeting S (no. 5, Corollary 2 to Proposition 11) and a fortiori \(x \notin p\) ; this establishes the first assertion. To prove the second, it is sufficient to prove

LEMMA 2. Every prime ideal of a ring A contains a minimal prime ideal of A.

It is sufficient, by Zorn's Lemma, to show that the set P of prime ideals, ordered by the relation \(\supset\) , is inductive. Now, if G is a non-empty totally ordered subset of P, the intersection \(p_{0}\) of the ideals \(p \in G\) is also a prime ideal: for if \(x \notin p_{0}\) and \(y \notin p_{0}\) , there exists an ideal \(p \in G\) such that \(x \notin p\) and \(y \notin p\) , whence \(xy \notin p\) and a fortiori \(xy \notin p_{0}\) .

Remark. In § 4, no. 3, Corollary 3 to Proposition 14, we shall show that in a Noetherian ring the set of minimal prime ideals is finite; moreover we shall see later that every decreasing sequence of prime ideals in a Noetherian ring is stationary.

COROLLARY 1. The radical of an ideal a in a ring A is the intersection of the prime ideals containing a and it is also the intersection of the minimal elements of this set of prime ideals.

COROLLARY 2. For an ideal a of a ring A we denote by r(a) the radical of a. Then, for two ideals of a, b of A,

\[
\mathfrak {r} (a \cap b) = \mathfrak {r} (a b) = \mathfrak {r} (a) \cap \mathfrak {r} (b);
\]

in particular, if \(\mathfrak{a} \subset \mathfrak{b}\) , then \(\mathfrak{r}(\mathfrak{a}) \subset \mathfrak{r}(\mathfrak{b})\) .

For a prime ideal to contain \(\mathfrak{a} \cap \mathfrak{b}\) (or \(\mathfrak{ab}\) ), it is necessary and sufficient that it contains one of the ideals \(\mathfrak{a}, \mathfrak{b}\) ( \(\S 1\) , no. 1, Proposition 1).

PROPOSITION 14. For two ideals a, b of a ring to be relatively prime, it is necessary and sufficient that their radicals r(a) and r(b) be so.

The necessity of the condition is obvious since \(a \subset \mathfrak{r}(a)\) and \(b \subset \mathfrak{r}(b)\) ; the condition is sufficient by §1, no. 2, Proposition 3.

PROPOSITION 15. In a ring A let \(\mathfrak{a}\) be an ideal and \(\mathfrak{b}\) a finitely generated ideal contained in the radical of \(\mathfrak{a}\) . Then there exists an integer \(k > 0\) such that \(\mathfrak{b}^k \subset \mathfrak{a}\) .

Let \((b_{i})_{1\leqslant i\leqslant n}\) be a system of generators of b. By hypothesis there exists an integer h such that \(b_{i}^{h}\in\mathfrak{a}\) for \(1\leqslant i\leqslant n\) . If a product of nh elements, each of which is a linear combination of the \(b_{i}\) with coefficients in A, is expanded, each term is a multiple of a product of nh factors, each of which is equal to a \(b_{i}\) ; among these factors, at least h have the same index i for some i, hence, the product belongs to a and nh is the required integer k.

COROLLARY. In a Noetherian ring the nilradical is a nilpotent ideal.

PROPOSITION 16. Let B be a ring and A a subring of B. For every minimal prime ideal p of A there exists a minimal prime ideal q of B such that q ∩ A = p.

We write \(S = A - p\) ; then the ring \(A_p = S^{-1}A\) is identified with a subring of \(S^{-1}B\) (no. 1, Proposition 2) and on the other hand \(A_p\) has only a single prime ideal \(p'\) since \(p\) is minimal (no. 5, Proposition 11). As \(S^{-1}B\) is not reduced to 0 (since it contains \(A_p\) ), it has at least one prime ideal \(r'\) and therefore \(r' \cap A_p = p'\) ; if \(j = i_B^S\) and we write \(r = j^{-1}(r')\) , then

\[
i _ {\mathbf {A}} ^ {\mathbf {S}} (\mathfrak {r} \cap \mathbf {A}) \subset \mathfrak {r} ^ {\prime} \cap \mathbf {A} _ {p} = p ^ {\prime}
\]

hence \(\mathfrak{r} \cap \mathbf{A} \subset \mathfrak{p}\) and, as \(\mathfrak{p}\) is minimal, \(\mathfrak{r} \cap \mathbf{A} = \mathfrak{p}\) ; moreover, \(\mathfrak{r}\) is prime in \(\mathbf{B}\) ; if \(\mathfrak{q}\) is a minimal prime ideal of \(\mathbf{B}\) contained in \(\mathfrak{r}\) (Lemma 1), then a fortiori \(\mathfrak{q} \cap \mathbf{A} = \mathfrak{p}\) since \(\mathfrak{p}\) is minimal.

DEFINITION 5. A ring A is called reduced if its nilradical is reduced to 0, in other words if no element \(\neq 0\) of A is nilpotent.

If N is the nilradical of a ring A, A/N is reduced, for if the class mod. N of an element \(x \in A\) is nilpotent in A/N, this means that \(x^{h} \in N\) for some integer h, hence \(x^{hk} = 0\) for some integer k and \(x \in N\) .

PROPOSITION 17. Let A be a ring and \(\mathfrak{A}\) its nilradical. For every multiplicative subset S of A, \(S^{-1}\mathfrak{A}\) is the nilradical of \(S^{-1}A\) . In particular, if A is reduced, then \(S^{-1}A\) is reduced.

If \(x \in \mathbf{A}\) , \(s \in \mathbf{S}\) satisfy \((x / s)^n = x^n / s^n = 0\) , there exists \(s' \in \mathbf{S}\) such that \(s' x^n = 0\) (no. 1, Remark 3) and a fortiori \((s' x)^n = 0\) , hence \(s' x \in \mathfrak{N}\) and \(x / s = s' x / s' s \in \mathbf{S}^{-1} \mathfrak{N}\) ; the converse is immediate.

\subsection{Modules of fractions of tensor products and homomorphism modules}

PROPOSITION 18. Let A be a ring and S a multiplicative subset of A.

\begin{enumerate}
\def\labelenumi{(\roman{enumi})}
\tightlist
\item
  If M and N are two A-modules, the \(S^{-1}A\) -modules ( \(S^{-1}M\) ) \(\otimes_{A}N\) , \(M\otimes_{A}(S^{-1}N)\) , ( \(S^{-1}M\) ) \(\otimes_{S^{-1}A}(S^{-1}N)\) and \(S^{-1}(M\otimes_{A}N)\) are canonically isomorphic.
\item
  If \(M'\) and \(N'\) are two \(S^{-1}A\) -modules, the canonical homomorphism

\[
\mathbf {M} ^ {\prime} \otimes_ {\mathbf {A}} \mathbf {N} ^ {\prime} \rightarrow \mathbf {M} ^ {\prime} \otimes_ {\mathbf {S} ^ {- 1} \mathbf {A}} \mathbf {N} ^ {\prime}
\]

derived from the A-bilinear mapping \((x', y') \mapsto x' \otimes y'\) from \(M' \times N'\) to \(M' \otimes_{S^{-1}} A N'\) is bijective.

\end{enumerate}
Assertion (i) is an immediate consequence of the definition \(\mathbf{S}^{-1}\mathbf{M} =\)

\(M \otimes_{A} S^{-1} A\) and the associativity of tensor products, which gives to within canonical isomorphisms

\[
\begin{array}{r l} \left(\mathrm{S} ^ {- 1} \mathrm{M} \otimes_ {\mathrm{S} ^ {- 1} \mathrm{A}} \left(\mathrm{S} ^ {- 1} \mathrm{N}\right) = (\mathrm{S} ^ {- 1} \mathrm{M}) \otimes_ {\mathrm{S} ^ {- 1} \mathrm{A}} \left(\mathrm{S} ^ {- 1} \mathrm{A} \otimes \mathrm{N}\right) = (\mathrm{S} ^ {- 1} \mathrm{M}) \otimes_ {\mathrm{A}} \mathrm{N} \right. \\ = (\mathrm{S} ^ {- 1} \mathrm{A}) \otimes_ {\mathrm{A}} \mathrm{M} \otimes_ {\mathrm{A}} \mathrm{N} = \mathrm{S} ^ {- 1} (\mathrm{M} \otimes_ {\mathrm{A}} \mathrm{N}). \end{array}
\]

To prove (ii), we note first that in \(\mathbf{M}'\) and \(\mathbf{N}'\) , considered as A-modules, the homotheties induced by the elements \(s \in S\) are bijective, hence \(\mathbf{M}' = \mathbf{S}^{-1}\mathbf{M}'\) and \(\mathbf{N}' = \mathbf{S}^{-1}\mathbf{N}'\) (no. 2, Remark 4) and similarly \(\mathbf{S}^{-1}(\mathbf{M}' \otimes_{\mathbf{A}} \mathbf{N}') = \mathbf{M}' \otimes_{\mathbf{A}} \mathbf{N}'\) ; (ii) is then a special case of (i).

COROLLARY. Let M be an A-module and a an ideal of A. The sub- \(\mathbf{S}^{-1}\mathbf{A}\) -modules \((\mathbf{S}^{-1}\mathfrak{a})(\mathbf{S}^{-1}\mathbf{M})\) , \(\mathfrak{a}(\mathbf{S}^{-1}\mathbf{M})\) , \((\mathbf{S}^{-1}\mathfrak{a})j(\mathbf{M})\) (where \(j: \mathbf{M} \to \mathbf{S}^{-1}\mathbf{M}\) is the canonical mapping) and \(\mathbf{S}^{-1}(\mathfrak{a}\mathbf{M})\) of \(\mathbf{S}^{-1}\mathbf{M}\) are identical. In particular, if \(\mathfrak{a}\) and \(\mathfrak{b}\) are two ideals of A, then

\[
(S ^ {- 1} a) (S ^ {- 1} b) = a (S ^ {- 1} b) = (S ^ {- 1} a) b = S ^ {- 1} (a b).
\]

Remark. Let M, N, P be three A-modules, \(f: M \times N \to P\) an A-bilinear mapping and \(S^{-1}f: (S^{-1}M) \times (S^{-1}N) \to (S^{-1}P)\) the \(S^{-1}A\) -bilinear mapping obtained from f by extending the ring of scalars to \(S^{-1}A\) (Algebra, Chapter IX, § 1, no. 4, Proposition 1). Let \(i: M \to S^{-1}M\) , \(j: N \to S^{-1}N\) be the canonical homomorphisms; it follows immediately that, if Q is the sub-A-module of P generated by \(f(M \times N)\) , \(S^{-1}Q\) is the sub- \(S^{-1}A\) -module of \(S^{-1}P\) generated by \((S^{-1}f)(i(M) \times j(N))\) .

PROPOSITION 19. Let A be a ring and S a multiplicative subset of A.

\begin{enumerate}
\def\labelenumi{(\roman{enumi})}
\tightlist
\item
  If M and N are two A-modules and M is finitely generated (resp. finitely presented), the canonical homomorphism (Chapter I, § 2, no. 10, formula (10))

\[
\mathrm{S} ^ {- 1} \operatorname{Hom} _ {\mathrm{A}} (\mathrm{M}, \mathrm{N}) \rightarrow \operatorname{Hom} _ {\mathrm{S} ^ {- 1} \mathrm{A}} (\mathrm{S} ^ {- 1} \mathrm{M}, \mathrm{S} ^ {- 1} \mathrm{N})
\]

is injective (resp. bijective).

\item
  If \(M'\) , \(N'\) are two \(S^{-1}A\) -modules, the canonical bijection

\[
\operatorname{Hom} _ {\mathbf {S} ^ {- 1} \mathbf {A}} (\mathbf {M} ^ {\prime}, \mathbf {N} ^ {\prime}) \rightarrow \operatorname{Hom} _ {\mathbf {A}} (\mathbf {M} ^ {\prime}, \mathbf {N} ^ {\prime})
\]

is bijective.

\end{enumerate}
As \(S^{-1}A\) is a flat A-module, (i) is a particular case of Chapter I, §2, no. 10, Proposition 11. On the other hand, we have already remarked that every A-homomorphism of \(S^{-1}A\) -modules is necessarily an \(S^{-1}A\) -homomorphism, in the course of proving Proposition 3 of no. 2; whence (ii).

PROPOSITION 20. Let A, A' be two rings, \(\rho: \mathbf{A} \to \mathbf{A}'\) a homomorphism, S a multiplicative subset of A, \(S' = \rho(S)\) and \(\rho': S^{-1}A \to S'^{-1}A'\) the homomorphism corresponding to \(\rho\) (no. 1, Proposition 2).

\begin{enumerate}
\def\labelenumi{(\roman{enumi})}
\tightlist
\item
  For every \(\mathbf{A}'\) -module \(\mathbf{M}'\) there exists a unique \(\mathbf{S}^{-1}\mathbf{A}\) -isomorphism

\[
j \colon \mathrm{S} ^ {- 1} \rho_ {*} (\mathrm{M} ^ {\prime}) \rightarrow \rho_ {*} ^ {\prime} (\mathrm{S} ^ {- 1} \mathrm{M} ^ {\prime})
\]

such that \(j(m' / s) = m' / \rho(s)\) for all \(m' \in \mathbf{M}'\) , \(s \in \mathbf{S}\) .

\item
  For every A-module M, there exists a unique isomorphism

\[
j ^ {\prime} \colon (\mathrm{S} ^ {- 1} \mathrm{M}) \otimes_ {\mathrm{S} ^ {- 1} \mathrm{A}} (\mathrm{S} ^ {\prime - 1} \mathrm{A} ^ {\prime}) \rightarrow \mathrm{S} ^ {\prime - 1} (\mathrm{M} \otimes_ {\mathrm{A}} \mathrm{A} ^ {\prime})
\]

of \(\mathbf{S}^{\prime -1}\mathbf{A}^{\prime}\) -modules such that \(j^{\prime}((m / s)\otimes (a^{\prime} / s^{\prime})) = (m\otimes a^{\prime}) / (\rho (s)s^{\prime})\)

\end{enumerate}
\begin{enumerate}
\def\labelenumi{(\roman{enumi})}
\item
  If we consider \(S^{\prime-1}M^{\prime}\) as an A-module by means of the composite homomorphism \(i_{M^{\prime}}^{S^{\prime}}\circ\rho\) , the homotheties induced by the elements of S are bijective, hence there exists a unique homomorphism j with the stated property (no. 2, Proposition 3). As \(\rho(S)=S^{\prime}\) , j is surjective; moreover, if \(m^{\prime}\in M^{\prime}\) , \(s\in S\) , \(m^{\prime}/\rho(s)=0\) , there exists \(t^{\prime}\in S^{\prime}\) such that \(t^{\prime}m^{\prime}=0\) ; as there exists \(t\in S\) such that \(\rho(t)=t^{\prime}\) , \(t^{\prime}.m^{\prime}=0\) in \(\rho_{*}(M^{\prime})\) , hence \(m^{\prime}/s=0\) in \(S^{-1}\rho_{*}(M^{\prime})\) .
\item
  As \((\mathbf{S}^{-1}\mathbf{M})\otimes_{\mathbf{S}^{-1}\mathbf{A}}(\mathbf{S}^{\prime -1}\mathbf{A}^{\prime}) = (\mathbf{M}\otimes_{\mathbf{A}}\mathbf{S}^{-1}\mathbf{A})\otimes_{\mathbf{S}^{-1}\mathbf{A}}(\mathbf{S}^{\prime -1}\mathbf{A}^{\prime})\) and

\[
\mathbf {S} ^ {- 1} (\mathbf {M} \otimes_ {\mathbf {A}} \mathbf {A} ^ {\prime}) = (\mathbf {M} \otimes_ {\mathbf {A}} \mathbf {A} ^ {\prime}) \otimes_ {\mathbf {A} ^ {\prime}} (\mathbf {S} ^ {- 1} \mathbf {A} ^ {\prime}),
\]

the existence of \(j'\) follows from the associativity of tensor products.

\end{enumerate}
\subsection{Application to algebras}

Let A be a ring, B an A-algebra (not necessarily associative or commutative and not necessarily possessing a unit element) and S a multiplicative subset of A. We know that it is possible to define on the \(S^{-1}A\) -module \(S^{-1}B = B \otimes_{A} S^{-1}A\) a canonical \(S^{-1}A\) -algebra structure, said to be obtained by extension of the ring of scalars to \(S^{-1}A\) (Algebra, Chapter III, § 3) and under which the product \((x/s)(y/t)\) is equal to \((xy)/(st)\) . If e is the unit element of B, e/1 is the unit element of \(S^{-1}B\) and if B is associative (resp. commutative), so is \(S^{-1}B\) .

Let A be a ring and M an A-module; we denote by T(M) (resp. ∧(M), S(M)) the tensor algebra (resp. exterior algebra, symmetric algebra) of M (Algebra, Chapter III). It is known that for every commutative A-algebra C there exists a unique isomorphism \(j\) of \(\mathsf{T}(\mathbf{M}) \otimes_{\mathbf{A}} \mathbf{C}\) onto \(\mathsf{T}(\mathbf{M} \otimes_{\mathbf{A}} \mathbf{C})\) (resp. of \(\bigwedge (\mathbf{M}) \oplus_{\mathbf{A}} \mathbf{C}\) onto \(\bigwedge (\mathbf{M} \otimes_{\mathbf{A}} \mathbf{C})\) , of \(\mathsf{S}(\mathbf{M}) \otimes_{\mathbf{A}} \mathbf{C}\) onto \(\mathsf{S}(\mathbf{M} \otimes_{\mathbf{A}} \mathbf{C})\) ) such that

\[
i (x \otimes 1) = x \otimes 1
\]

for all \(x \in \mathbf{M}\) , \(\mathbf{M}\) being canonically identified with a submodule of \(T(\mathbf{M})\) (resp. \(\bigwedge (\mathbf{M}), S(\mathbf{M})\) ) (loc. cit.). Then it is seen in particular that for every multiplicative subset \(\mathbf{R}\) of \(\mathbf{A}\) there are canonical isomorphisms

\[
\begin{array}{c} \mathbf {R} ^ {- 1} \mathsf {T} (\mathbf {M}) \to \mathsf {T} (\mathbf {R} ^ {- 1} \mathbf {M}), \qquad \mathbf {R} ^ {- 1} \bigwedge (\mathbf {M}) \to \bigwedge (\mathbf {R} ^ {- 1} \mathbf {M}), \\ \mathbf {R} ^ {- 1} \mathsf {S} (\mathbf {M}) \to \mathsf {S} (\mathbf {R} ^ {- 1} \mathbf {M}) \end{array}
\]

which reduce to the identity on \(R^{-1}M\) .

\subsection{Modules of fractions of graded modules}

Let A be a graded ring, M a graded A-module and \(\Delta\) the degree monoid; we shall assume in this no. that \(\Delta\) is a group. Recall (Algebra, Chapter II, § 11) that A and M are respectively direct sums of additive groups

\[
\mathbf {A} = \bigoplus_ {i \in \Delta} \mathbf {A} _ {i}, \quad \mathbf {M} = \bigoplus_ {i \in \Delta} \mathbf {M} _ {i}
\]

with \(\mathbf{A}_i\mathbf{A}_j\subset \mathbf{A}_{i + j}\) and \(\mathrm{A}_i\mathrm{M}_j\subset \mathrm{M}_{i + j}\) for all \(i,j\in \Delta\) . Let S be a multiplicative subset of A all of whose elements are homogeneous. For all \(i\in \Delta\) , we write \(\mathbf{S}_i = \mathbf{S}\cap \mathbf{A}_i\) and denote by \((\mathrm{S}^{-1}\mathbf{M})_i\) the set of elements \(m^{\prime}\) of \(\mathbf{S}^{-1}\mathbf{M}\) for which there exist elements \(j\) , \(k\) of \(\Delta\) , an element \(m\in \mathbf{M}_j\) and an element \(s\in \mathbf{S}_k\) such that \(j - k = i\) and \(m^{\prime} = m / s\) . If \((m_q')_{1\leqslant q\leqslant r}\) is a finite family of elements of \(\mathbf{S}^{-1}\mathbf{M}\) such that

\[
m _ {q} ^ {\prime} \in (\mathrm{S} ^ {- 1} \mathrm{M}) _ {i (q)},
\]

there exist elements \(j(q) \in \Delta\) and \(k \in \Delta\) , elements \(m_q \in \mathrm{M}_{j(q)}\) ( \(1 \leqslant q \leqslant r\) ) and \(s \in \mathrm{S}_k\) such that \(m_q' = m_q / s\) for \(1 \leqslant q \leqslant r\) (no. 1, Remark 2).

PROPOSITION 21. The ring \(\mathbf{S}^{-1}\mathbf{A}\) with the family \(((\mathbf{S}^{-1}\mathbf{A})_i)\) is a graded ring and \(\mathbf{S}^{-1}\mathbf{M}\) with the family \(((\mathbf{S}^{-1}\mathbf{M})_i)\) is a graded module over the graded ring \(\mathbf{S}^{-1}\mathbf{A}\) . The canonical mappings \(i_{\mathbf{A}}^{\mathbf{S}}\) and \(i_{\mathbf{M}}^{\mathbf{S}}\) are homogeneous of degree 0.

Let \(m \in M_j\) , \(s \in S_k\) , \(m' \in M_{j'}\) , \(s' \in S_{k'}\) and suppose that \(j - k = j' - k' = i\) ; then \((m/s) - (m'/s') = (s'm - sm')/ss'\) and \(s'm - sm' \in M_{j+k'} = M_{j'+k}\) and \(ss' \in S_{k+k'}\) , hence \((m/s) - (m'/s') \in (\mathrm{S}^{-1}\mathrm{M})_i\) by definition; this shows that the \((\mathrm{S}^{-1}\mathrm{M})_i\) are additive subgroups of \(S^{-1}\mathrm{M}\) . The sum of these groups is the whole of \(S^{-1}\mathrm{M}\) : for every \(x \in S^{-1}\mathrm{M}\) may be written as \(m/s\) where \(m \in \mathbf{M}\) , \(s \in \mathbf{S}\) ; \(s\) is homogeneous by hypothesis and \(m\) is a sum of homogeneous elements \(m_j\) ; hence \(x\) is the sum of the \(m_j / s\) , each of which belongs to a subgroup \((\mathbf{S}^{-1}\mathbf{M})_i\) . Finally the sum of the \((\mathbf{S}^{-1}\mathbf{M})_i\) is direct; for let us consider a finite family of elements \(x_q\) ( \(1 \leqslant q \leqslant n\) ) such that \(x_q \in (\mathbf{S}^{-1}\mathbf{M})_{i(q)}\) , where the indices \(i(q)\) are distinct, and suppose that \(\sum_{q=1}^{n} x_q = 0\) . Each \(x_q\) may be written as \(x_q = m_q / s\) where \(s \in S_k\) and \(m_q \in M_{i(q)+k}\) ; the hypothesis implies that there exists \(s' \in S\) such that \(s'\left(\sum_{q=1}^{n} m_q\right) = 0\) ; if the \(s'm_q\) were not all zero, we would have a contradiction since, if \(s' \in S_d\) , then \(s'm_q \in M_{i(q)+d}\) and the \(i(q) + d\) are all distinct. It follows that \(x_q = 0\) for every index \(q\) .

It is immediately verified that, if \(a \in (\mathrm{S}^{-1}\mathrm{A})_i\) and \(x \in (\mathrm{S}^{-1}\mathrm{M})_j\) , then \(ax \in (\mathrm{S}^{-1}\mathrm{M})_{i+j}\) . Applying this result to the case \(\mathbf{M} = \mathbf{A}\) , we see first that \(\mathrm{S}^{-1}\mathbf{A}\) is a ring graded by the \((\mathrm{S}^{-1}\mathrm{A})_i\) ; then we see that \(\mathrm{S}^{-1}\mathrm{M}\) is graded module over \(\mathrm{S}^{-1}\mathrm{A}\) . Finally, as \(1 \in \mathrm{A}_0\) , \(i_{\mathrm{A}}^{\mathrm{S}}\) and \(i_{\mathrm{M}}^{\mathrm{S}}\) are homogeneous of degree 0.

PROPOSITION 22. Let A (resp. B) be a graded ring of type \(\Delta\) , M (resp. N) a graded module over the graded ring A (resp. B), S (resp. T) a multiplicative subset of A (resp. B) all of whose elements are homogeneous, \(f: \mathrm{A} \to \mathrm{B}\) a homogeneous homomorphism of degree 0 such that \(f(\mathrm{S}) \subset \mathrm{T}\) and \(u: \mathrm{M} \to \mathrm{N}\) an A-linear mapping which is homogeneous of degree \(k\) . Then the homomorphism \(f': \mathrm{S}^{-1}\mathrm{A} \to \mathrm{T}^{-1}\mathrm{B}\) derived from \(f\) (no. 1, Proposition 2) is homogeneous of degree 0 and the \((\mathrm{S}^{-1}\mathrm{A})\) -linear mapping

\[
u ^ {\prime} \colon \mathrm{S} ^ {- 1} \mathrm{M} \rightarrow \mathrm{T} ^ {- 1} \mathrm{N}
\]

derived from f and u (no. 2, Proposition 5) is homogeneous of degree k.

This follows immediately from the equations \(f'(a / s) = f(a) / f(s)\) and \(u'(m / s) = u(m) / f(s)\) .

Finally we note that, if E is a graded A-algebra and S a multiplicative subset of A consisting of homogeneous elements, \(S^{-1}E\) with its \((S^{-1}A)\) -algebra structure (no. 8) and the graduation \((S^{-1}E)_{t}\) is a graded \(S^{-1}A\) -algebra, as follows immediately from the definitions.

\section{LOCAL RINGS. PASSAGE FROM THE LOCAL TO THE GLOBAL}

\subsection{Local rings}

PROPOSITION 1. Let A be a ring and I the set of non-invertible elements of A. The set I is the union of the ideals of A which are distinct from A. Moreover, the following conditions are equivalent:

\begin{enumerate}
\def\labelenumi{(\alph{enumi})}
\item
  I is an ideal.
\item
  The set of ideals of \(\mathbf{A}\) distinct from \(\mathbf{A}\) has a greatest element.
\item
  A has a unique maximal ideal.
\end{enumerate}

The relation \(x \in I\) is equivalent to \(1 \notin xA\) and hence \(xA \neq A\) . If \(a\) is an ideal of \(A\) distinct from \(A\) and \(x \in a\) , then \(xA \subset a\) , hence \(xA \neq A\) and \(x \in I\) . Hence every ideal of \(A\) distinct from \(A\) is contained in \(I\) and every element \(x \in I\) belongs to a principal ideal \(xA \neq A\) . This proves the first assertion, which immediately implies the equivalence of (a), (b) and (c).

Remark (1). Note that, if (c) holds, I is the Jacobson radical of the ring A (Algebra, Chapter VIII, § 6, no. 3, Definition 3).

DEFINITION 1. A ring A is called a local ring if it satisfies the equivalent conditions (a), (b) and (c) of Proposition 1. The quotient of A by its Jacobson radical (which is then the unique maximal ideal of A) is called the residue field of A.

DEFINITION 2. Let A, B be two local rings and m, n their respective maximal ideals. A homomorphism \(u: \mathbf{A} \to \mathbf{B}\) is called local if \(u(\mathfrak{m}) \subset \mathfrak{n}\) .

This amounts to saying that \(\bar{u}^{1}(\mathfrak{n}) = \mathfrak{m}\) , for \(\bar{u}^{1}(\mathfrak{n})\) is then an ideal containing m and not containing l and hence equal to m. Taking quotients, we then derive canonically from u an injective homomorphism A/m → B/n from the residue field of A to that of B.

\textbf{Examples}\par

\begin{enumerate}
\def\labelenumi{(\arabic{enumi})}
\item
  A field is a local ring. A ring which is reduced to 0 is not a local ring.
\item
  Let A be a local ring and k its residue field. The ring of formal power series \(B = A[[X_{1}, \ldots, X_{n}]]\) is a local ring, for the non-invertible elements of B are the formal power series whose constant terms are not invertible in A (Algebra, Chapter IV, § 5, no. 6, Proposition 4). The canonical injection of A into B is a local homomorphism and the corresponding injection of residue fields is an isomorphism.
\item
  Let \(\mathfrak{b}\) be an ideal of a ring A which is only contained in a single maximal ideal m; then A/b is a local ring with maximal ideal m/b and residue field canonically isomorphic to A/m. This applies in particular to the case \(b = m^{k}\) , m being any maximal ideal of A ( \(\S\) 1, no. 1, Corollary to Proposition 1). If A itself is a local ring with maximal ideal m, then for every ideal \(b \neq A\) of A, A/b is a local ring, the canonical homomorphism \(A \to A/b\) a local homomorphism and the corresponding homomorphism of residue fields an isomorphism.
\item
  Let X be a topological space, \(x_{0}\) a point of X and A the ring of germs at the point \(x_{0}\) of real-valued functions continuous in a neighbourhood of \(x_{0}\) (General Topology, Chapter I, §6, no. 10). Clearly, for the germ at \(x_{0}\) of a continuous function f to be invertible in A, it is necessary and sufficient that \(f(x_{0}) \neq 0\) , as that implies that \(f(x) \neq 0\) in a neighbourhood of \(x_{0}\) . The ring A is therefore a local ring whose maximal ideal m is the set of germs of functions which are zero at \(x_{0}\) ; taking quotients, the mapping \(g \mapsto g(x_{0})\) of A to R gives an isomorphism of the residue field A/m onto R.
\end{enumerate}

PROPOSITION 2. Let A be a ring and p a prime ideal of A. The ring \(A_{p}\) is local; its maximal ideal is the ideal \(pA_{p} = p_{p}\) , generated by the canonical image of p in \(A_{p}\) ; its residue field is canonically isomorphic to the field of fractions of A/p.

Let \(S = A - p\) and \(j: A \to A_p\) be the canonical homomorphism; the hypothesis that \(p\) is prime implies that \(p\) is saturated with respect to \(S\) , hence \(j^{1}(\mathfrak{pA}_{\mathfrak{p}}) = \mathfrak{p}\) (§ 2, no. 4, Proposition 10) and, as the ideals of A not meeting S are those contained in \(\mathfrak{p}\) , the first two assertions are special cases of § 2, no. 5, Proposition 11 (ii). Moreover, if \(f\) is the canonical homomorphism \(A \to A/\mathfrak{p}\) , \(f(S)\) is the set of elements \(\neq 0\) of the integral domain \(A/\mathfrak{p}\) and hence the last assertion is a special case of § 2, no. 5, Proposition 11 (i).

DEFINITION 3. Let A be a ring and p a prime ideal of A. The ring \(A_{p}\) is called the local ring of A at p, or the local ring of p, when there is no ambiguity.

Remark (2). If A is a local ring and m its maximal ideal, the elements of A --- m are invertible (Proposition 1) and hence \(A_{m}\) is canonically identified with A (§ 2, no. 1, Remark 5).

\textbf{Examples}\par

\begin{enumerate}
\def\labelenumi{(\arabic{enumi})}
\setcounter{enumi}{4}
\tightlist
\item
  Let \(p\) be a prime number. The local ring \(\mathbf{Z}_{(p)}\) is the set of rational numbers \(a / b\) , where \(a, b\) are rational integers with \(b\) prime to \(p\) ; the residue field of \(\mathbf{Z}_{(p)}\) is isomorphic to the prime field \(\mathbf{F}_p = \mathbf{Z} / (p)\) .
\end{enumerate}

* (6) Let V be an affine algebraic variety, A the ring of regular functions on V, W an irreducible subvariety of V and p the (necessarily prime) ideal of A consisting of the functions which are zero at every point of W. The ring \(\mathbf{A}_{\mathfrak{p}}\) is called the local ring of W on V. *

PROPOSITION 3. Let A be a ring, p a prime ideal of A and S = A - p.~For every ideal \(\mathfrak{b}'\) of \(\mathbf{A}_{\mathfrak{p}}\) distinct from \(\mathbf{A}_{\mathfrak{p}}\) , let \(\mathfrak{b}\) be the ideal \((i_{\mathbf{A}}^{\mathbf{S}})^{-1}(\mathfrak{b}')\) of A so that \(\mathfrak{b}' = \mathfrak{bA}_{\mathfrak{p}}\) .

\begin{enumerate}
\def\labelenumi{(\roman{enumi})}
\item
  Let \(f\) be the canonical homomorphism \(\mathbf{A} \to \mathbf{A} / \mathfrak{b}\) . The homomorphism from \(\mathbf{A}_{\mathfrak{p}}\) to \((\mathbf{A} / \mathfrak{b})_{\mathfrak{p} / \mathfrak{b}}\) canonically associated with \(f\) (§ 2, no. 1, Proposition 2) is surjective and its kernel is \(\mathfrak{b}'\) , which defines, by taking quotients, a canonical isomorphism of \(\mathbf{A}_{\mathfrak{p}} / \mathfrak{b}'\) onto \((\mathbf{A} / \mathfrak{b})_{\mathfrak{p} / \mathfrak{b}}\) .
\item
  The mapping \(\mathfrak{b}' \to \mathfrak{b} = (i_{\mathbf{A}}^{\mathbb{S}})^{-1}(\mathfrak{b}')\) , restricted to the set of prime ideals of \(\mathbf{A}_{\mathfrak{p}}\) , is an isomorphism (with respect to inclusion) of this set onto the set of prime ideals of \(\mathbf{A}\) contained in \(\mathfrak{p}\) . If \(\mathfrak{b}'\) is prime in \(\mathbf{A}_{\mathfrak{p}}\) , there exists an isomorphism of the ring \(\mathbf{A}_{\mathfrak{b}}\) onto the ring \((\mathbf{A}_{\mathfrak{p}})_{\mathfrak{b}'}\) , which maps \(a / s\) to \((a / 1) / (s / 1)\) for all \(a \in \mathbf{A}\) , \(s \in \mathbf{A} - \mathfrak{b}\) ).
\end{enumerate}

This is just a particular case of § 2, no. 5, Proposition 11.

\textbf{Remarks}\par

\begin{enumerate}
\def\labelenumi{(\arabic{enumi})}
\setcounter{enumi}{2}
\item
  If \(\mathfrak{a}\) is an ideal of \(\mathbf{A}\) not contained in \(\mathfrak{p}\) , then \(\mathfrak{a}\mathrm{A}_{\mathfrak{p}} = \mathrm{A}_{\mathfrak{p}}\) and \((\mathrm{A} / \mathfrak{a})_{\mathfrak{p}} = 0\) (§ 2, no. 5, Remark).
\item
  Let A, B be two rings, \(\rho: A \to B\) a homomorphism, \(q\) a prime ideal of B and \(p\) the prime ideal \(\rho^{-1}(q)\) of A. As \(\rho(A - p) \subset B - q\) , a canonical homomorphism \(\rho_q: A_p \to B_q\) is derived from \(\rho\) ( \(\S 2\) , no. 1, Proposition 2) and it is immediate that \(\rho_q(pA_p) \subset qB_q\) , hence \(\rho_q\) is a local homomorphism.
\end{enumerate}
